% ------------------------------------------------------- V C dim
% chapter
% 
% 09/27/2010 16:56:36 Finished a pass. Would benefit form another
% iteration.  --------------------------------------------------------

\makeatletter
\def\input@path{{styles/}}
\makeatother
\ifx\bookMode\undefined%
    %-----------------------------------------------------------------------------
% Master Prefix File for Randomized Algorithms Class Notes & Book
% Unified self-contained configuration for LuaLaTeX and XeLaTeX with Memoir
%-----------------------------------------------------------------------------

\RequirePackage{iftex}
\ifPDFTeX
  \PackageError{engine-guard}{This document requires LuaLaTeX or XeLaTeX}{%
    Please recompile using LuaLaTeX or XeLaTeX instead of pdfLaTeX.%
  }
  \stop
\fi

% Always use memoir
\documentclass[12pt,oneside,fleqn]{memoir}

%-----------------------------------------------------------------------------
% Typography, Fonts & Math Packages
%-----------------------------------------------------------------------------
\usepackage{mathtools}
\usepackage{newcomputermodern}   % Loads fontspec & unicode-math

% OpenType math aliases and fixes
\providecommand{\square}{\mdlgwhtsquare}
\def\Box{\mdlgwhtsquare}
\def\Diamond{\mdlgwhtdiamond}
\def\blacksquare{\mdlgblksquare}
\def\leadsto{\rightsquigarrow}

\usepackage{alphabeta}
\usepackage[normalem]{ulem} % Provides \sout
\usepackage{mleftright}
\usepackage{microtype}
\usepackage{xspace}
\usepackage{xurl}
\usepackage[dvipsnames,table]{xcolor}
\usepackage{graphicx}
\usepackage{caption}
\usepackage{booktabs,array,colortbl,hhline,multirow,tabularx}
\usepackage[inline]{enumitem}
\usepackage[thicklines]{cancel}
\renewcommand\CancelColor{\color{purple}}
\usepackage{euscript}
\usepackage{marvosym}
\let\Cross\undefined
\let\Square\undefined
\usepackage{bbding}
\usepackage{stmaryrd}
\usepackage{epigraph}
\usepackage{makeidx}
\usepackage{perpage}
\usepackage{tikz}

%-----------------------------------------------------------------------------
% Boondox / mathcalb font declaration
%-----------------------------------------------------------------------------
\DeclareFontFamily{U}{BOONDOX-calo}{\skewchar\font=45 }
\DeclareFontShape{U}{BOONDOX-calo}{m}{n}{
  <-> s*[1.05] BOONDOX-r-calo}{}
\DeclareFontShape{U}{BOONDOX-calo}{b}{n}{
  <-> s*[1.05] BOONDOX-b-calo}{}
\DeclareMathAlphabet{\mathcalb}{U}{BOONDOX-calo}{m}{n}
\SetMathAlphabet{\mathcalb}{bold}{U}{BOONDOX-calo}{b}{n}
\DeclareMathAlphabet{\mathbcalb}{U}{BOONDOX-calo}{b}{n}

%-----------------------------------------------------------------------------
% Algorithm Environment (formerly salgorithm.sty)
%-----------------------------------------------------------------------------
\newcommand{\ProgWordX}[1]{{\small\textbf{#1}}{\ }}
\newcommand{\Do}{\ProgWordX{do}\xspace}
\newcommand{\End}{\ProgWordX{end}\xspace}
\newcommand{\Return}{\ProgWordX{return}\xspace}%
\newcommand{\Procbegin}{\ProgWordX{begin}}
\newcommand{\If}{\ProgWordX{if}\xspace}
\newcommand{\Then}{\ProgWordX{then}\xspace}
\newcommand{\Else}{\ProgWordX{else}\xspace}
\newcommand{\For}{\ProgWordX{for}\xspace}
\newcommand{\While}{{\small\bfseries while}\ }

\newcommand{\xbeginlgox}{\begin{minipage}{1in}\begin{tabbing}
           \quad\=\qquad\=\qquad\=\qquad\=\qquad\=\qquad\=\qquad\=\kill}
\newcommand{\xendlgox}{\end{tabbing}\end{minipage}}

\newenvironment{algorithmEnv}{%
   \begin{tabular}{|l|}\hline\xbeginlgox\bigskip}%
   {\xendlgox\\\hline\end{tabular}}

\newenvironment{program}{%
   \begin{minipage}{4.0in}%
   \begin{tabbing}%
       \ \ \ \ \=\ \ \ \=\ \ \ \ \=\ \ \ \ \=\ \ \ \ \=%
      \ \ \ \ \=\ \ \ \ \=\ \ \ \ \=\ \ \ \ \=%
      \ \ \ \ \=\ \ \ \ \=\ \ \ \ \=\ \ \ \ \=\kill%
}{%
   \end{tabbing}%
   \end{minipage}%
}

\newcommand{\AlgorithmBox}[1]{%
   \fbox{\begin{program} #1 \end{program}}%
}

\definecolor{AlgorithmColor}{cmyk}{0.07,0.90,0,0.34}
\providecommand{\Algorithm}[1]{{\AlgorithmI{#1}\index{algorithm!#1@{\AlgorithmI{#1}}}}}
\newcommand{\CodeComment}[1]{\textcolor{blue}{\texttt{#1}}}

%-----------------------------------------------------------------------------
% Path & Directory Resolution
%-----------------------------------------------------------------------------
\newcommand{\RootDIR}{../}
\IfFileExists{../../../book.tex}{\renewcommand{\RootDIR}{../../../}}{}
\IfFileExists{../../book.tex}{\renewcommand{\RootDIR}{../../}}{}
\IfFileExists{../book.tex}{\renewcommand{\RootDIR}{../}}{}
\IfFileExists{./book.tex}{\renewcommand{\RootDIR}{}}{}

\providecommand{\File}[1]{#1}
\providecommand{\RootFile}[1]{\RootDIR #1}

%-----------------------------------------------------------------------------
% QED Duck Symbol
%-----------------------------------------------------------------------------
\newlength{\CharHeight}
\AtBeginDocument{\setlength{\CharHeight}{\fontcharht\font`X}}

\IfFileExists{\RootDIR shared/qed_duck.pdf}{%
   \providecommand{\myqedsymbol}{%
      \reflectbox{\includegraphics[height=1.4\CharHeight]{\RootDIR shared/qed_duck.pdf}}%
   }%
}{%
   \IfFileExists{shared/qed_duck.pdf}{%
      \providecommand{\myqedsymbol}{%
         \reflectbox{\includegraphics[height=1.4\CharHeight]{shared/qed_duck.pdf}}%
      }%
   }{%
      \providecommand{\myqedsymbol}{$\blacksquare$}%
   }%
}

%-----------------------------------------------------------------------------
% BibLaTeX Setup (Pure Biber, No BibTeX)
%-----------------------------------------------------------------------------
\usepackage[bibencoding=utf8,style=alphabetic,refsegment=chapter,
indexing=cite,
backend=biber]{biblatex}

\ExecuteBibliographyOptions{%
   maxnames = 40,%
   maxalphanames = 4,%
   maxcitenames = 3,%
   minalphanames = 3,%
   minnames = 1,%
   giveninits = true,%
   doi = false,%
   url = false,%
   isbn = false%
}

\renewbibmacro{in:}{}

\AtEveryBibitem{%
   \ifentrytype{inproceedings}{%
      \clearname{editor}%
      \clearlist{editor}%
      \clearfield{editor}%
      \clearlist{publisher}%
      \clearfield{series}%
      \clearlist{location}%
      \clearfield{publisher}%
   }{}%
   \ifentrytype{book}{%
      \clearname{publisher}%
   }{}%
}
\AtEveryCitekey{%
   \ifentrytype{inproceedings}{%
      \clearlist{publisher}%
   }{}%
   \clearname{editor}%
   \clearlist{location}%
}
\DeclareFieldFormat[inproceedings]{date}{}
\DeclareFieldFormat[inproceedings]{pages}{#1, \printfield{year}}
\DeclareFieldFormat[inproceedings]{booktitle}{\mkbibemph{#1,\nopunct}}

\DeclareFieldFormat{journaltitle}{\mkbibemph{#1}\isdot}

\DeclareFieldFormat[article]{number}{\mkbibparens{#1}}
\DeclareFieldFormat[article]{date}{}
\DeclareFieldFormat[article]{pages}{#1, \printfield{year}}
\DeclareFieldFormat[article]{issue}{\nopunct#1}
\renewbibmacro*{issue+date}{%
   \iffieldundef{issue}{}%
   {%
      \printtext{\unspace}%
      (\printfield{issue})%
   }%
   \iffieldundef{pages}%
   {\addcomma\printfield{year}}%
   {}%
}

\renewcommand{\bibpagespunct}{%
   \ifentrytype{article}{\addcolon}{\addcomma}\space%
}

\DeclareFieldFormat[article]{titlecase}{\MakeSentenceCase{#1}}
\DeclareFieldFormat[inproceedings]{titlecase}{\MakeSentenceCase{#1}}
\DeclareFieldFormat[incollection]{titlecase}{\MakeSentenceCase{#1}}
\DeclareFieldFormat[book]{titlecase}{\MakeSentenceCase{#1}}

\renewbibmacro*{journal}{%
   \iffieldundef{journaltitle} {} {\printtext[journaltitle]{%
         \printfield[myplain]{journaltitle}%
         \setunit{\subtitlepunct}%
         \printfield[myplain]{journalsubtitle}}}}
\DeclareFieldFormat{myplain}{#1}

\renewbibmacro*{booktitle}{%
   \ifboolexpr{ test {\iffieldundef{booktitle}} and test
      {\iffieldundef{booksubtitle}} } {} {\printtext[booktitle]{%
         \printfield[myplain]{booktitle}%
         \setunit{\subtitlepunct}%
         \printfield[myplain]{booksubtitle}}%
      \newunit}%
   \printfield{booktitleaddon}}

\renewbibmacro*{volume+number+eid}{%
   \setunit*{\addcomma\addspace}%
   \printfield{volume}%
   \printfield{number}%
   \setunit{\addcomma\space}%
   \printfield{eid}%
}

\newbibmacro{string+doi}[1]{%
   \iffieldundef{doi}{%
      \iffieldundef{url}{%
         #1%
      }{%
         \href{\thefield{url}}{#1}%
      }%
   }{%
      \href{http://dx.doi.org/\thefield{doi}}{#1}}%
}
\DeclareFieldFormat*{title}{\usebibmacro{string+doi}{\mkbibemph{#1}}}
\renewcommand*{\multicitedelim}{\addcomma\space}
\DeclareFieldFormat[inproceedings]{title}{\usebibmacro{string+doi}{{#1}}}
\DeclareFieldFormat[article]{title}{\usebibmacro{string+doi}{{#1}}}
\DeclareFieldFormat[incollection]{title}{\usebibmacro{string+doi}{{#1}}}
\DeclareFieldFormat[book]{titlecase}{{#1}}

\DeclareCiteCommand{\tcite}%
{\usebibmacro{prenote}}%
{\usebibmacro{citeindex}\usebibmacro{cite}}%
{\multicitedelim}%
{\usebibmacro{postnote}}

%-----------------------------------------------------------------------------
% Hyperref & Visual Formatting
%-----------------------------------------------------------------------------
\usepackage{hyperref}
\hypersetup{%
   breaklinks,%
   colorlinks=true,%
   linkcolor=[rgb]{0.45,0.0,0.0},%
   citecolor=[rgb]{0,0,0.45}%
}

\definecolor{DeepNavy}{HTML}{002366}
\definecolor{RoyalBlue}{HTML}{004E92}
\definecolor{SteelBlue}{HTML}{4682B4}
\definecolor{TealBlue}{HTML}{008080}
\definecolor{DarkSlate}{HTML}{2F4F4F}
\definecolor{DustyGrey}{HTML}{696969}

\makeatletter
\AtBeginDocument{\let\@elt\relax}
\makeatother

\headstyles{bringhurst}
\chapterstyle{dash}

\renewcommand{\chapnamefont}{\normalfont\huge\bfseries\color{DeepNavy}}
\renewcommand{\chapnumfont}{\normalfont\huge\bfseries\color{DeepNavy}}
\renewcommand{\chaptitlefont}{\normalfont\Huge\bfseries\color{DeepNavy}}

\counterwithin{footnote}{chapter}

\setsubparaheadstyle{\normalfont\normalsize\itshape\color{DustyGrey}}
\setlength{\beforechapskip}{-20pt} 
\setlength{\afterchapskip}{30pt}
\setsecnumformat{\csname the#1\endcsname.\quad}

%-----------------------------------------------------------------------------
% Standalone Chapter vs Master Book Modes
%-----------------------------------------------------------------------------
\newcommand{\Copyright}{This work is licensed under the Creative Commons Attribution-Noncommercial 3.0 License. To view a copy of this license, visit \url{http://creativecommons.org/licenses/by-nc/3.0/} or send a letter to Creative Commons, 171 Second Street, Suite 300, San Francisco, California, 94105, USA.}

\newcommand{\RealChapterEndInner}{%
   {%
      \begingroup%
      \emergencystretch=1em%
      \printbibliography[heading=subbibliography,segment=\therefsegment]{}%
      \endgroup%
   }%
}

\providecommand{\BookMode}[1]{}
\providecommand{\StandAloneMode}[1]{#1}

\ifx\bookMode\undefined
   % STANDALONE CHAPTER MODE
   \renewcommand{\StandAloneMode}[1]{#1}
   \renewcommand{\BookMode}[1]{}

   \newcommand{\ChapterNumPage}[2]{%
      \if\relax\detokenize{#1}\relax
         \setcounter{chapter}{0}%
      \else\ifnum#1>0
         \setcounter{chapter}{\numexpr#1-1\relax}%
      \else
         \setcounter{chapter}{0}%
      \fi\fi
      \if\relax\detokenize{#2}\relax
         \setcounter{page}{1}%
      \else\ifnum#2>0
         \setcounter{page}{#2}%
      \else
         \setcounter{page}{1}%
      \fi\fi
   }

   \DeclareRobustCommand{\setheadingcolor}{\color{RoyalBlue}}
   \setsecheadstyle{\setheadingcolor\normalfont\Large\bfseries}
   \setsubsecheadstyle{\normalfont\large\bfseries\color{SteelBlue}}
   \setsubsubsecheadstyle{\normalfont\normalsize\bfseries\color{TealBlue}}
   \setparaheadstyle{\normalfont\normalsize\bfseries\color{DarkSlate}}

   \DeclareDocumentCommand{\Chapter}{s o m}{%
      \IfBooleanTF{#1}{%
         \IfValueTF{#2}{\chapter*[]{#2}}{\chapter*[]{#3}}%
      }{%
         \IfValueTF{#2}{\chapter*[]{#2}}{\chapter[]{#3}}%
      }%
      \setcounter{footnote}{0}%
   }
   \newcommand{\ChapterPrefix}{}
   \newcommand{\ChapterGen}{\pagestyle{plain}\vspace*{-0.7cm}}
   \newcommand{\ChapterEnd}{%
      \RealChapterEndInner{}%
      \end{document}%
   }

   \renewcommand{\File}[1]{#1}
   \renewcommand{\RootFile}[1]{../#1}
\else
   % MASTER BOOK MODE
   \renewcommand{\StandAloneMode}[1]{}
   \renewcommand{\BookMode}[1]{#1}

   \newcommand{\ChapterNumPage}[2]{}

   \DeclareDocumentCommand{\Chapter}{s o m}{%
      \IfBooleanTF{#1}{%
         \IfValueTF{#2}{\chapter*[#2]{#3}}{\chapter*{#3}}%
         \ifx\Dir\undefined\else
            \typeout{CHAPTER_MAP: \Dir/\CurrentChapFile => chap: 0, page: \thepage}%
         \fi
      }{%
         \IfValueTF{#2}{\chapter[#2]{#3}}{\chapter{#3}}%
         \ifx\Dir\undefined\else
            \typeout{CHAPTER_MAP: \Dir/\CurrentChapFile => chap: \thechapter, page: \thepage}%
         \fi
      }%
   }
   \newcommand{\ChapterPrefix}{%
      \vspace{-1cm}%
      \subsection*{\small
         {598 - Class notes for Randomized Algorithms}\\
         {Sariel Har-Peled}\\
         {\today}}%
   }
   \newcommand{\ChapterGen}{}
   \newcommand{\ChapterEnd}{\RealChapterEndInner{}}

   \renewcommand{\File}[1]{\Dir/#1}
   \renewcommand{\RootFile}[1]{#1}

   \newcommand{\InputChap}[2]{%
      \def\Dir{#1}%
      \def\CurrentChapFile{#2}%
      \typeout{}%
      \typeout{------------------------------------------------------------}%
      \typeout{PROCESSING FILE: #1/#2}%
      \typeout{}%
      \input{#1/#2}%
   }
\fi

%-----------------------------------------------------------------------------
% Fragments & Quotes
%-----------------------------------------------------------------------------
\newcommand{\InputExt}[1]{%
   \IfFileExists{#1}%
   {\input{#1}}%
   {\IfFileExists{../#1}%
      {\input{../#1}}%
      {\IfFileExists{../../#1}%
         {\input{../../#1}}%
         {\typeout{ERROR: Unable to find #1}}%
      }%
   }%
}
\newcommand{\FragmentName}[1]{\RootFile{fragment/#1.tex}}

\newwrite\FragOutHandle
\newwrite\FragProbeOut
\newread\FragProbeIn
\newif\ifInJunkDir
\InJunkDirfalse

% Detect whether XeTeX was invoked with -output-directory=junk
\edef\FragProbeSecret{\number\uniformdeviate 1000000}
\immediate\openout\FragProbeOut=frag_probe.tmp
\immediate\write\FragProbeOut{\FragProbeSecret}
\immediate\closeout\FragProbeOut

\openin\FragProbeIn=junk/frag_probe.tmp
\ifeof\FragProbeIn
  \InJunkDirfalse
\else
  \read\FragProbeIn to \FragReadSecret
  \closein\FragProbeIn
  \def\trimspace#1 {#1}
  \edef\FragReadSecretClean{\expandafter\trimspace\FragReadSecret}
  \ifx\FragReadSecretClean\FragProbeSecret
    \InJunkDirtrue
  \else
    \InJunkDirfalse
  \fi
\fi

\newcommand{\FragWritePrefix}{%
  \ifx\bookMode\undefined
    % Standalone chapter mode
    \ifInJunkDir ../../fragment/\else ../fragment/\fi
  \else
    % Master book mode
    \ifInJunkDir ../fragment/\else fragment/\fi
  \fi
}

\DeclareDocumentEnvironment{fragment}{m +b}{%
   \immediate\openout\FragOutHandle=\FragWritePrefix#1.tex
   \immediate\write\FragOutHandle{\unexpanded{#2}}%
   \immediate\closeout\FragOutHandle
   #2%
}{}

\makeatletter
\DeclareDocumentCommand{\IncFragment}{m}{%
   \begingroup
      \let\label\@gobble
      \InputExt{\FragmentName{#1}}%
   \endgroup
}
\makeatother

\newcommand{\QuoteOpen}{}
\newcommand{\QuoteClose}{}
\newcommand{\QuoteWrite}[2]{}

\newcommand{\Quote}[2]{%
   \SaveIndent
   \QuoteWrite{#1}{#2}%
   \begin{center}%
      {\footnotesize
         \begin{quote}
             \RestoreIndent
             {#1}\\
             \hspace*{2cm}-- {#2.}
         \end{quote}%
      }%
      \index{quotation!#2}%
   \end{center}%
   \bigskip
}

\newcommand{\QuotePExt}[3]{%
   \hfill%
   \begin{minipage}{0.95\linewidth}%
       \footnotesize {#1}%
       \rule{\linewidth}{0.3mm}%
       \hfill #2, #3%
   \end{minipage}%
}

\epigraphfontsize{\small\itshape}
\setlength{\epigraphwidth}{0.6\textwidth}
\newcommand{\QuoteExt}[3]{%
   \vspace*{-1.6cm}%
   \epigraph{{\footnotesize #1}}{{#2}, {#3}}%
   \vspace*{-1.2cm}%
}

\providecommand{\delX}[1]{{\color{red}\sout{#1}}}
\providecommand{\newX}[1]{{\color{blue}#1}}
\providecommand{\chgY}[2]{{\color{red}\sout{#1}}{\color{blue}#2}}
\providecommand{\remX}[1]{{\color{purple}\textbf{[Remark: #1]}}}
\newcommand{\SpellIgnore}[1]{#1}
\newcommand{\EQ}{\mathrel{\;\;=\;\;}}

%-----------------------------------------------------------------------------
% Save a symbol that we know is going to get redefined.
\def\savesymbol#1{%
  \expandafter\let\expandafter\origsym\expandafter=\csname#1\endcsname
  \expandafter\let\csname orig#1\endcsname=\origsym
  \expandafter\let\csname#1\endcsname=\relax
}

% Restore a previously saved symbol, and rename the current one.
\def\restoresymbol#1#2{%
  \expandafter\let\expandafter\newsym\expandafter=\csname#2\endcsname
  \expandafter\global\expandafter\let\csname#1#2\endcsname=\newsym
  \expandafter\let\expandafter\origsym\expandafter=\csname orig#2\endcsname
  \expandafter\global\expandafter\let\csname#2\endcsname=\origsym
}

\usepackage{pictex}%
\savesymbol{triangle}%
\usepackage{tree}%
\restoresymbol{TreePKG}{triangle}%

% Definitions absorbed from sbook.sty
\providecommand{\MdPrc}{\nu}

\providecommand{\TwoFiguresBase}[4]{%
   \begin{tabular}{cc}
      \begin{minipage}{0.48\linewidth}
          \begin{center}
              {{{#1}}}
          \end{center}
      \end{minipage}
      &
      \begin{minipage}{0.48\linewidth}
          \begin{center}
              {{#3}}
          \end{center}
      \end{minipage}
      \\
      {#2} & {#4}
   \end{tabular}%
}

\providecommand{\TwoFigures}[6]{%
   \begin{tabular}{cc}
      \begin{minipage}{0.48\linewidth}
          \begin{center}
              {\IncludeGraphics[#1]{{#2}}}
          \end{center}
      \end{minipage}
      &
      \begin{minipage}{0.48\linewidth}
          \centerline{\IncludeGraphics[#4]{{#5}}}
      \end{minipage}
      \\
      {#3} & {#6}
   \end{tabular}%
}

\providecommand{\ThreeFigures}[6]{%
   \begin{tabular}{ccc}
        \begin{minipage}{0.30\linewidth}
            \IncludeGraphics[width=5.1cm]{{#1}}\\
        \end{minipage}
        &
        \begin{minipage}{0.30\linewidth}
            \IncludeGraphics[width=5.1cm]{{#2}}\\
        \end{minipage}
        &
        \begin{minipage}{0.30\linewidth}
            \IncludeGraphics[width=5.1cm]{{#3}}\\
        \end{minipage}
        \vspace{-1cm}
        \\
        {#4} & {#5}& {#6}
   \end{tabular}%
}

%-----------------------------------------------------------------------------
\providecommand{\remove}[1]{}
\providecommand{\NotAMSMode}[1]{#1}
\providecommand{\Solution}[1]{}
\providecommand{\Version}[1]{\hfill{Version: #1}}
\providecommand{\Distor}{\mathrm{dist}}
\providecommand{\Hardness}[1]{(h:#1)}
\providecommand{\QuoteOptExt}[3]{\QuoteExt{#1}{#2}{#3}}
\providecommand{\Hypercube}{\EuScript{H}}
\providecommand{\emphiC}[1]{\emphic{#1}{#1}}
\providecommand{\CoverTime}{\EuScript{C}}
\providecommand{\Nbr}{\Gamma}
\providecommand{\CommuteTime}[2]{\mathbf{CT}_{#1#2}}
\providecommand{\Resistance}{\mathbf{R}}
\providecommand{\EffectResist}[2]{\mathbf{R}_{#1#2}}
\providecommand{\NetworkE}{\EuScript{N}}
\providecommand{\trinom}[4]{\binom{#1}{#2 \;\; #3 \;\; #4}}
\providecommand{\Ranges}{\EuScript{R}}
\providecommand{\CClass}[1]{\ensuremath{\mathbf{#1}}}
\providecommand{\coCClass}[1]{\ensuremath{\mathrm{co-}}\CClass{#1}}
\providecommand{\CX}{\EuScript{C}}
\providecommand{\packet}{\mathbf{c}}
\providecommand{\packetA}{\mathbf{d}}
\providecommand{\mopt}{m_{\mathrm{opt}}}
\providecommand{\girth}{\mathop{\mathrm{girth}}}
\providecommand{\CR}{\mathop{\mathrm{cr}}}
\providecommand{\TRUE}{{\textcolor[named]{RedViolet}{\textsf{TRUE}}}\xspace}
\providecommand{\FALSE}{{\textcolor[named]{RedViolet}{\textsf{FALSE}}}\xspace}
\providecommand{\Part}[1]{\part{#1}}
\providecommand{\Section}[1]{\section{#1}}
\providecommand{\SubSection}[1]{\subsection{#1}}
\providecommand{\SubSubSection}[1]{\subsubsection{#1}}
\providecommand{\Exercise}[5]{%
\begin{exercise_h}[#1]%
   \exelab{#3}{}%
   #4%
\end{exercise_h}%
}
\providecommand{\si}[1]{#1}
\providecommand{\myParagraph}[1]{\noindent\textbf{#1}}
\providecommand{\DefName}[1]{#1}
\providecommand{\deflabel}[1]{def_#1}
\providecommand{\Description}[1]{}

%-----------------------------------------------------------------------------
% Mathematical Notations & Environments (from notations.sty)
%-----------------------------------------------------------------------------
%-----------------------------------------------------------------------------
% List environments (questions & compactenum)
%-----------------------------------------------------------------------------

\newlist{compactenumA}{enumerate}{5}%
\setlist[compactenumA]{topsep=0pt,itemsep=-1ex,partopsep=1ex,parsep=1ex,%
   label=(\Alph*)}%

\newlist{compactenumI}{enumerate}{5}%
\setlist[compactenumI]{topsep=0pt,itemsep=-1ex,partopsep=1ex,parsep=1ex,%
   label=(\Roman*)}%

\newlist{compactenuma}{enumerate}{5}%
\setlist[compactenuma]{topsep=0pt,itemsep=-1ex,partopsep=1ex,parsep=1ex,%
   label=(\alph*)}%

\newlist{compactitem}{enumerate}{5}%
\setlist[compactitem]{topsep=0pt,itemsep=-1ex,partopsep=1ex,parsep=1ex,%
   label=\textbullet}%

\newlist{compactenumi}{enumerate}{5}%
\setlist[compactenumi]{topsep=0pt,itemsep=-1ex,partopsep=1ex,parsep=1ex,%
   label=(\roman*)}%

\setcounter{secnumdepth}{5}

\newcommand{\HST}{\TermI{HST}\xspace}
\newcommand{\Term}[1]{\textsf{#1}}
\newcommand{\TermI}[1]{\Term{#1}\index{#1@\Term{#1}}}

\DeclareMathAlphabet{\mathpzc}{OT1}{pzc}{m}{it}

\newcommand{\seg}{\mathsf{s}}

\newcommand{\Diam}{\mathrm{diam}}
\newcommand{\diam}{\mathrm{diam}}

\newcommand{\Copt}{C_{\mathrm{opt}}}
\newcommand{\CenterPnt}{\mathrm{cen}}

\newcommand{\parent}{\overline{\mathrm{p}}}
\newcommand{\parentX}[2][\!]{\parent\pth[#1]{#2}}

\newcommand{\FMS}{\Mh{\EuScript{X}}\index{metric space}}  
\newcommand{\FMSB}{\Mh{\EuScript{Y}}\index{metric space}}  
\newcommand{\FastCut}{\Algorithm{FastCut}\xspace}
\newcommand{\MinCut}{\Algorithm{MinCut}\xspace}
\newcommand{\MinCutRep}{\Algorithm{MinCutRep}\xspace}
\newcommand{\Contract}{\Algorithm{Contract}\xspace}

\newcommand{\Grid}{\mathrm{G}}
\newcommand{\ClosestPair}{\mathcal{CP}}
\newcommand{\SphereChar}{\ensuremath{\mathbb{S}}}
\newcommand{\Sphere}[1]{{\SphereChar^{(#1)}}}

\newcommand{\PntSet}{{\Mh{\mathsf{P}}}}%
\providecommand{\PntSetA}{{\Mh{\mathsf{Q}}}}%

\newcommand{\query}{\mathcalb{q}}

\newcommand{\pnt}{\mathsf{p}}
\newcommand{\pntA}{\mathsf{u}}
\newcommand{\pntB}{\mathsf{v}}
\newcommand{\pntC}{{\mathsf{x}}}
\newcommand{\pntD}{{\mathsf{y}}}
\newcommand{\pntE}{{\mathsf{z}}}

\newcommand{\pntCoord}{\Mh{p}}
\newcommand{\pntCoordA}{\Mh{q}}

%%%%%%%%%%%%%%%%%%%%%%%%%%%%%%%%%%%%%%%%%%%%%%%%%%%%%%%%%%%%%%%%%%%%% 
%%%%%%%%%%%%%%%%%%%%%%%%%%%%%%%%%%%%%%%%%%%%%%%%%%%%%%%%%%%%%%%%%%%%%
% Markov chain notations.

\newcommand{\MCSpc}{\Mh{\mathtt{S}}}
\newcommand{\MCPrb}{\Mh{\mathtt{P}}}
\newcommand{\MCPrbY}[2]{\MCPrb_{#1#2}}
\newcommand{\MCPrbM}{\Mh{\mathbf{P}}}%
\newcommand{\MCPrbT}[3]{\MCPrb_{#1#2}^{(#3)}}
\newcommand{\MCFPrbT}[3]{\Mh{\mathtt{r}}_{#1#2}^{(#3)}}
\newcommand{\MCVstProb}[2]{\Mh{\mathtt{f}}_{#1#2}}
\newcommand{\MCExSteps}[2]{\Mh{\mathtt{h}}_{#1#2}}
\newcommand{\MCPVec}[1]{\mathbf{q}^{(#1)}}
%%%%%%%%%%%%%%%%%%%%%%%%%%%%%%%%%%%%%%%%%%%%%%%%%%%%%%%%%%%%%%%%%%%%%

\newcommand{\Spread}{\Phi}
\newcommand{\VC}{\ensuremath{\mathsf{VC}}\xspace}
\newcommand{\MC}{\ensuremath{\mathsf{MC}}\xspace}
\newcommand{\hyperplane}{h}

\newcommand{\QuickSort}{\Algorithm{QuickSort}\xspace}
\newcommand{\Alg}{\Algorithm{Alg}\xspace}

\newcommand{\Graph}{\Mh{\mathsf{G}}}%
\newcommand{\GraphB}{\mathsf{H}}
\newcommand{\GraphC}{\mathsf{K}}
\newcommand{\GraphD}{\mathsf{F}}

\providecommand{\G}{}
\renewcommand{\G}{\Graph}

\newcommand{\RLP}{\ComplexityClass{RLP}\xspace}
\newcommand{\USTCON}{\ComplexityClass{USTCON}}
\newcommand{\STCON}{\overrightarrow{\ComplexityClass{STCON}}}

\newcommand{\Mtrx}{\mathsf{M}}
\newcommand{\TMtrx}{\mathsf{Q}}
\newcommand{\IMtrx}{\mathcal{I}}

\newcommand{\Expander}[3]{$\left[#1,#2,#3\right]$-expander}

\newcommand{\ChartX}[1]{\ensuremath{\chi_{#1}}}
\newcommand{\Field}{\mathbb{F}}
\newcommand{\LSpace}{\mathbb{G}}
\newcommand{\LD}[2]{\mathsf{L{D}}\pth{#1,#2}}
\newcommand{\SEVN}{\widehat{\lambda_2}}
\newcommand{\NX}[2]{#1\!\left[#2\right]}
\newcommand{\Entry}[2]{#1\!\left[#2\right]}
\newcommand{\NXExt}[3]{#1_{#2}\!\left[#3\right]}
\newcommand{\RangeI}[1]{\llbracket #1 \rrbracket}
\newcommand{\DialX}[1]{\gamma\pth{#1}}
\newcommand{\BiDialX}[1]{{\gamma}_2\pth{#1}}

\newcommand{\EV}[1]{\widehat{\lambda_{#1}}}
\newcommand{\TEV}[1]{\widehat{\lambda_{#1}}}
\newcommand{\EVal}{\widehat{\lambda}}
\newcommand{\Basis}{\mathcal{B}}

\newcommand{\ZProd}[2]{\ensuremath{#1 \, \text{\textcircled{z}} #2}}
\newcommand{\RProd}[2]{\ensuremath{#1 \, \text{\textcircled{r}} #2}}

\newcommand{\OneVec}[1]{\mathbf{1}^{#1}}
\newcommand{\DstrbV}[1]{\mathbf{\mathsf{q}}^{(#1)}}
\newcommand{\DstrbE}[2]{\mathsf{q}^{(#1)}_{#2}}

\newcommand{\DEdges}{\widetilde{\Edges}}

\newcommand{\VV}{\Mh{\mathsf{V}}}%
\newcommand{\EE}{\Mh{\mathsf{E}}}%

\newcommand{\Vertices}{\VV}%
\newcommand{\Edges}{\EE}%

\newcommand{\EdgesX}[1]{\Edges\pth{#1}}
\newcommand{\VerticesX}[1]{V\pth{#1}}

\newcommand{\StringX}[1]{\mathsf{s}\pth{#1}}

\newcommand{\BMtrx}{\mathsf{B}}
\newcommand{\WMtrx}{\mathsf{W}}
\newcommand{\NWMtrx}{\overline{\mathsf{W}}}
\newcommand{\Sequence}{\EuScript{S}}

\newcommand{\med}{\mathrm{med}}
\newcommand{\medX}[1]{\Mh{\mathrm{med}}\pth{#1}}

\newcommand{\medianC}{\Mh{\nu}}%

\newlength{\savedparindent}
\newcommand{\SaveIndent}{\setlength{\savedparindent}{\parindent}}
\newcommand{\RestoreIndent}{\setlength{\parindent}{\savedparindent}}

\newcommand{\NNN}{\mathbb{N}}%
\providecommand{\NN}{\mathbb{N}}%
\renewcommand{\Re}{\mathbb{R}}%
\newcommand{\reals}{\Re}%
\newcommand{\ZZ}{\mathbb{Z}}%

\newcommand{\level}{\mathop{\mathrm{level}}}

\providecommand{\TPDF}[2]{\texorpdfstring{#1}{#2}}

\newcommand{\Hint}[1]{\emph{[#1]}}
\newcommand{\Arr}{\mathop{\mathrm{\EuScript{A}}}}

\newcommand{\Array}{\mathcal{T}}%

\newcommand{\dirEdge}[2]{\left({{#1} \rightarrow {#2}} \right)}

%%%%%%%%%%%%%%%%%%%%%%%%%%%%%%%%%%%%%%%%%%%%%%%%%%%%%%%%%%%%%%%%%%%%%%
% Emphasize stuff...
\definecolor{blue25}{rgb}{0, 0, 11}
\newcommand{\emphic}[2]{%
   \textcolor{blue25}{%
      \bfseries{\emph{#1}}}%
   \index{#2}}
\newcommand{\emphi}[1]{\emphic{#1}{#1}}
\newcommand{\emphOnly}[1]{\textcolor{blue25}{\bfseries{\emph{#1}}}}

\definecolor{almostblack}{rgb}{0, 0, 0.3}
\providecommand{\emphw}[1]{{\textcolor{almostblack}{\emph{#1}}}}%

%
%%%%%%%%%%%%%%%%%%%%%%%%%%%%%%%%%%%%%%%%%%%%%%%%%%%%%%%%%%%%%%%%%%%%%%

\newcommand{\volX}[1]{\mathrm{vol}\pth{#1}}

\newcommand{\Interval}{\EuScript{I}}
\newcommand{\IntervalA}{\EuScript{I}}

\newcommand{\Exp}{\mathrm{Exp}}

\let\Ball\undefined
\newcommand{\Ball}{\mathbf{b}}
\newcommand{\hdim}{d}
\newcommand{\Domain}{\mathcal{D}}
\newcommand{\points}[1]{[\bfseries{#1 points}]}
\newcommand{\Event}{\Mh{\EuScript{E}}}%
\newcommand{\EventB}{\Mh{\EuScript{B}}}%
\newcommand{\EventF}{\Mh{\EuScript{F}}}%
\newcommand{\EventG}{\Mh{\EuScript{G}}}%

\newcommand{\distCmd}[1]{\left\| {#1}  \right\|}
\newcommand{\distY}[2]{\distCmd{#1 - #2}}

\newcommand{\NormalC}{\mathbf{N}}
\newcommand{\NormalY}[2]{\NormalC\pth{#1, #2}}%
\newcommand{\NormalDistOneDim}{\NormalC}
\newcommand{\NormalDistX}[1]{\NormalC^{#1}}
\DeclareMathOperator{\lca}{lca}
\newcommand{\RangeSpace}{\Mh{\mathsf{S}}}%
\newcommand{\RangeSpaceA}{\mathsf{T}}

\newcommand{\range}{\mathbf{r}}
\newcommand{\rangeA}{\mathbf{r}'}

\newcommand{\BadProb}{\varphi}

\newcommand{\EntropyChar}{\mathbb{H}}
\newcommand{\Entropy}[1]{\EntropyChar\pth{#1}}
\newcommand{\EntropyX}[1]{\EntropyChar\pth[]{#1}}

\newcommand{\Matrix}{\mathsf{M}}
\newcommand{\MatrixB}{\mathsf{B}}
\newcommand{\MatrixC}{\mathsf{C}}
\newcommand{\MatrixD}{\mathsf{D}}

\renewcommand{\th}{th\xspace}

\newcommand{\Family}{\mathcal{F}}
\newcommand{\FamilyB}{\mathcal{G}}
\newcommand{\FamilyC}{\mathcal{H}}
\newcommand{\FMSC}{\EuScript{Y}\index{metric space}}

\newcommand{\median}{\overline{\mathsf{m}}}
\newcommand{\QuickSelect}{\Algorithm{Quick{}Select}\xspace}

\providecommand{\row}{\mathbf{r}}%
\newcommand{\rowc}{{r}}%

\newcommand{\vecb}{\mathbf{b}}

\newcommand{\vecA}{\mathbf{v}}
\newcommand{\vecB}{\mathbf{u}}
\newcommand{\vecRand}{\mathbf{r}}

\newcommand{\vecH}{\mathbf{h}}
\newcommand{\vecV}{\mathbf{v}}

\newcommand{\vecX}{\mathbf{x}}
\newcommand{\vecY}{\mathbf{y}}

\newcommand{\DistBinom}[2]{\mathrm{Bin}\pth{#1,#2}}

\newcommand{\DistGeom}[1]{\mathrm{Geom}\pth{#1}}

\newcommand{\Sample}{\mathsf{R}}

\newcommand{\lcm}{\mathop{\mathrm{lcm}}}
\newcommand{\bdiv}{\mathop{\mathrm{div}}}%
\newcommand{\divides}{\mid}%
\newcommand{\poly}{\mathrm{poly}}%

\newcommand{\Group}{\mathcal{G}} 
\newcommand{\GroupA}{\mathcal{H}} 

\newcommand{\algGCD}{\Algorithm{EuclidGCD}\xspace}%

\newcommand{\gOp}{\times}%
\newcommand{\gIdent}{\mathsf{i}}%

\newcommand{\ordX}[2][\!]{\mathop{\mathrm{o{r}d}}\pth[#1]{#2}}

\newcommand{\lSym}[2]{%
   \left(
       \ds {#1}\divides{#2}%
   \right) %
}
\newcommand{\jSym}[2]{\left\llbracket%
       \ds {#1}\divides{#2}%
       \right\rrbracket%
}
\newcommand{\Line}{\ell}

\newcommand{\LevelChar}{\mathbb{L}}
\newcommand{\AtMostKLevelChar}{\LevelChar_{\leq k}}
\newcommand{\AtMostKLevel}[1]{\LevelChar_{\leq k} \pth{ #1}}

\newcommand{\LineSet}{\mathsf{L}}%
\newcommand{\ArrX}[1]{\EuScript{A}\pth{#1}}%

\providecommand{\Matousek}{Matou{\v s}ek\xspace}
\providecommand{\Erdos}{Erd{\H o}s\xspace}
\providecommand{\Lovasz}{Lov\'asz\xspace}

\newcommand{\ProblemC}[1]{{\textsf{\textcolor{red}{#1}\index{problem!#1}}}}
\DefineNamedColor{named}{OliveGreen} {cmyk}{0.64,0,0.95,0.40}
\providecommand{\ComplexityClass}[1]{{{\textcolor[named]{OliveGreen}{%
            \textsc{#1}}}}}
\providecommand{\NPHard}{{\ComplexityClass{NP-Hard}}%
   \index{NP!hard}\xspace}
\providecommand{\NPComplete}{\ComplexityClass{NP-Complete}%
   \index{NP!complete}\xspace}
\providecommand{\APXHard}{\ComplexityClass{APX-Hard}%
   \index{APX!Hard}\xspace}
\providecommand{\NPC}{\ComplexityClass{NPC}%
   \index{NP!complete}\xspace}
\providecommand{\NP}{\ComplexityClass{NP}%
   \index{NP}\xspace}
\providecommand{\POLYT}{\ComplexityClass{P}\xspace}

\newcommand{\envFrame}[1]{\begin{center}\fbox{\begin{minipage}[c]{0.75\textwidth}{#1}\end{minipage}}\end{center}}

\newcommand{\mysqrt}[1]{\ts\!\!\sqrt{#1}}
\newcommand{\SegSet}{\mathsf{S}}
\newcommand{\ArrVD}[1]{\Arr^{|} \pth{#1}}
\newcommand{\RIC}{\textsf{RIC}\index{random incremental
      construction}\xspace}
\newcommand{\ArrB}{\EuScript{B}}

\newcommand{\TreeOpt}{\EuScript{T}_\mathrm{opt}}
\newcommand{\TreeOptA}{\EuScript{T}_\mathrm{opt}'}

\newcommand{\DepthExt}[2]{\mathrm{depth}_{#1}\pth{#2}}
\newcommand{\depthX}[1]{\mathrm{depth}\pth{#1}}

\newcommand{\TreeA}{\EuScript{U}}
\newcommand{\TreeB}{\EuScript{V}}

\newcommand{\Trap}{\sigma}
\newcommand{\TrapA}{\sigma'}
\newcommand{\TrapB}{\sigma''}

\newcommand{\SampleA}{\mathsf{R}'}

\newcommand{\AllRegions}{\mathcal{T}}
\newcommand{\RegionsSet}[1]{\mathcal{F}\pth{#1}}
\newcommand{\RegionsHeavy}[2]{\mathcal{F}_{\!\geq#1}\pth{#2}}
\newcommand{\NumOfRegionsZero}[1]{\ExChar \!f\pth{#1}}
\newcommand{\NumOfRegions}[2]{\ExChar\! f_{\!\geq #1}\pth[]{#2}}
\newcommand{\ObjSet}{\mathsf{S}}

\newcommand{\WeightX}[2][\!]{\omega \pth[#1]{#2}}
\newcommand{\Edge}{\mathsf{e}}

\newcommand{\Face}{\mathsf{f}}

\newcommand{\DefSet}[1]{D\pth{#1}}
\newcommand{\KillSet}[1]{K\pth{#1}}

\newcommand{\CSProbFull}[4]{
   \,\rho_{\!\!\stackrel{\rule[-.005cm]{0cm}{0.1cm}}{#1,
       #4}}\pth{\MakeSBig #2,#3}}

\newcommand{\disk}{D}
\newcommand{\diskA}{D'}

\newcommand{\IntervalSet}{\widehat{\mathcal{I}}}

\newcommand{\GroundSet}{\textsf{X}}
\newcommand{\FGroundSet}{\mathsf{x}}
\newcommand{\RangeSet}{{\mathcal{R}}}
\newcommand{\RangeSetA}{\mathcal{R}'}
\newcommand{\MeasureChar}{\overline{m}}
\newcommand{\Measure}[1]{\MeasureChar\pth{#1}}
\newcommand{\sMeasure}[1]{\overline{s}\pth{#1}}

\newcommand{\SetA}{\Mh{S}}%
\newcommand{\SetB}{\Mh{T}}%
\newcommand{\SetC}{\Mh{N}}%

\newcommand{\VCProj}[2]{#1_{|#2}}
\newcommand{\VCDim}{\mathrm{dim}_{\VC}}
\newcommand{\Dim}{\delta}
\newcommand{\DimA}{\delta'}
\newcommand{\PntSetB}{{{\mathsf{Q}}}}
\newcommand{\PntSetC}{{{\mathsf{R}}}}
\newcommand{\Growth}[2]{{\EuScript{G}}_{#1}\pth{#2}}
\newcommand{\ShatterFunc}[1]{\pi_{#1}}
\newcommand{\Distribution}{\EuScript{D}}
\newcommand{\Error}{\EuScript{E}}

\newcommand{\bd}{{\partial}}%
\newcommand{\eps}{{\varepsilon}}%

\newcommand{\clmlab}[1]{\label{claim:#1}}
\newcommand{\clmref}[1]{\HLink{Claim}{claim:#1}}

\newcommand{\CHX}[1]{{\mathcal{CH}}\pth{#1}}

\newcommand{\ts}{\hspace{0.6pt}}

\def\qedsymbol{\myqedsymbol}%

\newcommand{\qedsign}{\hfill{\hfill\myqedsymbol}}
\newcommand{\qedsignBig}{\hfill{\hfill\rule{4mm}{4mm}}}
\newcommand{\QED}{{\hfill{\hfill\qedsymbol}}}

\newcommand{\num}{\alpha}

\newcommand{\NEW}[1]{#1}

\newcommand{\ThreeFiguresExt}[9]{%
   \begin{tabular}{ccc}
     \begin{minipage}{0.31\linewidth}
         \begin{center}
             {\IncludeGraphics[#1]{{#2}}}
         \end{center}
     \end{minipage}
     &
       \begin{minipage}{0.31\linewidth}
           \centerline{\IncludeGraphics[#4]{{#5}}}
       \end{minipage}
     &
       \begin{minipage}{0.31\linewidth}
           \centerline{\IncludeGraphics[#7]{{#8}}}
       \end{minipage}
     \\
     ~\\[-0.2cm]
     {#3} & {#6} & {#9}
   \end{tabular}}

\newcommand{\IncludeGraphics}[2][]{%
   \typeout{}%
   \typeout{Graphics: #2}%
   \typeout{\ includegraphics[#1]{#2}}%
   \includegraphics[#1]{#2}
   \typeout{}%
}

%    -----------------------------------------------------------------
%    -----------------------------------------------------------------
%    -----------------------------------------------------------------
%    Handling references
%    -----------------------------------------------------------------
%    -----------------------------------------------------------------
%    -----------------------------------------------------------------
\newcommand{\HLinkShort}[2]{\hyperref[#2]{#1\ref*{#2}}}
\newcommand{\HLink}[2]{\hyperref[#2]{#1~\ref*{#2}}}
\newcommand{\HLinkPage}[2]{\hyperref[#2]{#1~\ref*{#2}%
      $_\text{p\pageref{#2}}$}}

\newcommand{\HLinkSuffix}[3]{\hyperref[#2]{#1\ref*{#2}{#3}}}
\newcommand{\HLinkPageSuffix}[3]{\hyperref[#2]{#1\ref*{#2}%
      #3$_\text{p\pageref{#2}}$}}

\newcommand{\figlab}[1]{\label{fig:#1}}
\newcommand{\figref}[1]{\HLink{Figure}{fig:#1}}

\newcommand{\exelab}[1]{\label{exer:#1}}
\newcommand{\exeref}[1]{Exercise~\ref{exer:#1}}

\providecommand{\deflab}[1]{\label{def:#1}}
\newcommand{\defref}[1]{\HLink{Definition}{def:#1}}
\newcommand{\defrefY}[2]{\hyperref[def:#2]{#1}}
\newcommand{\defrefpage}[1]{\HLinkPage{Definition}{def:#1}}

\newcommand{\etal}{\textit{et~al.}\xspace}

\newcommand{\obslab}[1]{\label{observation:#1}}
\newcommand{\obsref}[1]{\HLink{Observation}{observation:#1}}
\newcommand{\obsrefY}[2]{\hyperref[observation:#2]{#1}}

\newcommand{\exmlab}[1]{\label{example:#1}}
\newcommand{\exmref}[1]{\HLink{Example}{example:#1}}

\newcommand{\remlab}[1]{\label{rem:#1}}
\newcommand{\remref}[1]{\HLink{Remark}{rem:#1}}%

\newcommand{\chaplab}[1]{\label{chap:#1}}

\newcommand{\itemlab}[1]{\label{item:#1}}
\newcommand{\itemref}[1]{\HLinkSuffix{}{item:#1}{}}

\newcommand{\tedlab}[1]{\label{tedium:#1}}
\newcommand{\tedref}[1]{\HLink{Tedium}{tedium:#1}}%

\newcommand{\factlab}[1]{\label{fact:#1}}%

\newcommand{\alglab}[1]{\label{algorithm:#1}}

\newcommand{\net}{\Mh{\mathcal{N}}}%
      
\newcommand{\lemlab}[1]{\label{lemma:#1}}
\newcommand{\lemref}[1]{\HLink{Lemma}{lemma:#1}}%
\newcommand{\lemrefY}[2]{\hyperref[lemma:#2]{#1}}
\newcommand{\lemrefpage}[1]{\HLinkPage{Lemma}{lemma:#1}}
\newcommand{\lemrefshort}[1]{\HLinkShort{L}{lemma:#1}}

\newcommand{\eqlab}[1]{\label{equation:#1}}
\newcommand{\Eqrefbare}[1]{\HLinkSuffix{(}{equation:#1}{)}}
\newcommand{\Eqref}[1]{\HLinkSuffix{Eq.~(}{equation:#1}{)}}
\newcommand{\Eqrefpage}[1]{\HLinkPageSuffix{Eq.~(}{equation:#1}{)}}

\newcommand{\seclab}[1]{\label{sec:#1}}
\newcommand{\secref}[1]{\HLink{Section}{sec:#1}}
\newcommand{\secrefY}[2]{\hyperref[sec:#2]{#1}}
\newcommand{\secrefpage}[1]{\HLinkPage{Section}{sec:#1}}

\newcommand{\thmlab}[1]{{\label{theo:#1}}}
\newcommand{\thmref}[1]{\HLink{Theorem}{theo:#1}}
\newcommand{\thmrefY}[2]{\hyperref[theo:#2]{#1}}
\newcommand{\thmrefpage}[1]{\HLinkPage{Theorem}{theo:#1}}

\newcommand{\corlab}[1]{\label{cor:#1}}
\newcommand{\corref}[1]{\HLink{Corollary}{cor:#1}}%

%    -----------------------------------------------------------------
%    -----------------------------------------------------------------
%    -----------------------------------------------------------------
%    Handling references
%    -----------------------------------------------------------------
%    -----------------------------------------------------------------
%    -----------------------------------------------------------------

\newcommand{\ceil}[1]{\left\lceil {#1} \right\rceil}
\newcommand{\floor}[1]{\left\lfloor {#1} \right\rfloor}
\newcommand{\brc}[1]{\left\{ {#1} \right\}}
\newcommand{\pth}[2][\!]{\mleft({#2}\mright)}%
\newcommand{\brk}[1]{\mleft[{#1}\mright]}%
\newcommand{\pbrc}[2][\!\!]{#1\left[ {#2} \bigr. \right]}
\newcommand{\pbrcS}[1]{\left[ {#1} \right]}

\newcommand{\dotprod}[2]{{#1}\cdot{#2}}
\newcommand{\DotProd}[2]{\permut{{#1},{#2}}}
\newcommand{\permut}[1]{\left\langle {#1} \right\rangle}

\newcommand{\MakeBig}{}
\newcommand{\MakesBig}{}
\newcommand{\MakeSBig}{}
\newcommand{\MakeVBig}{}

\newcommand{\DistChar}{\Mh{\mathsf{d}}}%
\newcommand{\DistSetY}[2]{\DistChar\pth{#1,#2}}
\newcommand{\distG}[2]{\DistChar_\Graph\pth{#1, #2}}
\newcommand{\DistMtrCharX}[1]{\DistChar_{#1}}

\newcommand{\dY}[2]{\DistChar\pth{#1, #2}}%
\newcommand{\dGY}[2]{\DistChar_\Graph\pth{#1, #2}}%

\newcommand{\cset}{\mathcal{C}}%

\newcommand{\MTR}{\mathcal{M}}
\newcommand{\DistMtr}[2]{\DistChar\pth{#1,#2}}

\newcommand{\dist}[1]{\left\| {#1} \right\|}
\newcommand{\norm}[1]{\left\| {#1} \right\|}

\providecommand{\ds}{\displaystyle}
\newcommand{\sep}[1]{\,\left|\, {#1} \bigr.\right.}%
\newcommand{\sepw}[1]{\left|\, {#1} \right.}%

\newcommand{\ProbExt}[2]{\mathop{\mathbf{Pr}}_{#1}\pbrcx{#2}}
\newcommand{\ProbLTR}{\mathbb{P}}%

\newcommand{\Prob}[1]{\mathop{\ProbLTR} \mleft[ #1 \mright]}%
\newcommand{\ProbCond}[2]{\mathop{\ProbLTR}\!\left[%
       #1 \;\middle\vert\; #2 \right]}
\newcommand{\ProbY}[2]{\ProbLTR^{}_{\!#1}\pbrc{#2}}

\newcommand{\ExChar}{\mathbb{E}}%
\newcommand{\ExSym}{\mathop{\ExChar}}%
\newcommand{\Ex}[2][\!]{\ExSym#1\pbrcx{#2}}
\newcommand{\ExExt}[2]{\ExSym_{#1}\!\pbrcx{#2}}
\newcommand{\ExExtS}[2]{\ExSym_{\substack{#1}}\!\pbrcx{#2}}
\newcommand{\ExCond}[2]{\ExSym\!\left[%
       #1 \;\middle\vert\; #2 \right]}

\newcommand{\VarSymb}{\mathop{\mathbb{V}}}
\newcommand{\Var}[1]{\mathop{\mathbb{V}}\!\pbrcx{#1}}

\newcommand{\pbrcx}[1]{\left[ {#1} \right]}%
\newcommand{\nfrac}[2]{{#1}/{#2}}
\newcommand{\cardin}[1]{\left| {#1} \right|}%

\newcommand\CCancel[2][black]{{\renewcommand\CancelColor{\color{#1}}\cancel{#2}}}

\newcommand{\IncGraphPage}[4][]{%
   \IncludeGraphics[page=#4,#1]{\File{#2/#3}}%
}

\newenvironment{enumEx}{
   \begin{compactenumA}
    }{\end{compactenumA}}

\newcommand{\cl}{{\mathrm{cl}}}

\usepackage[amsmath,thmmarks]{ntheorem}
\theoremseparator{.}%

\newtheorem{theorem:unnumbered}{Theorem}[section]
\newtheorem{theorem}{Theorem}[section]

\newtheorem{observation}[theorem]{Observation}
\newtheorem{conjecture}[theorem]{Conjecture}
\newtheorem{lemma}[theorem]{Lemma}
\newtheorem{corollary}[theorem]{Corollary}
\newtheorem{claim}[theorem]{Claim}
\newtheorem{assumption}[theorem]{Assumption}
\newtheorem{tedium}[theorem]{Tedium}

\theoremstyle{plain}%
\theoremheaderfont{\sffamily}
\theorembodyfont{\upshape}%

\newtheorem{defn}[theorem]{Definition}
\newtheorem{definition}[theorem]{Definition}

\newtheorem*{remark_unnumbered}[theorem:unnumbered]{Remark}%
\newtheorem*{remarks}[theorem]{Remarks}%
\newtheorem{remark}[theorem]{Remark}%
\newtheorem{fact}[theorem]{Fact}%
\newtheorem{example}[theorem]{Example}
\newtheorem{exercise}[theorem]{Exercise}
\newtheorem{problem}[theorem]{Problem}
\newtheorem{algorithm}[theorem]{Algorithm}%
\newtheorem{xca}[theorem]{Exercise}
\newtheorem{openproblem}[theorem]{Open Problem}
\newtheorem{exercise_h}[theorem]{Exercise}

% Proof environments
\theoremheaderfont{\em}%
\theorembodyfont{\upshape}%
\theoremstyle{nonumberplain}
\theoremseparator{}
\theoremsymbol{\myqedsymbol}
\newtheorem{proof}{Proof:}
\newtheorem{proofextin}{Proof}
\newtheorem{altproof}{Alternative proof:}

\theoremsymbol{}
\newtheorem{proofnoqed}{Proof:}

\theoremsymbol{\myqedsymbol}
\theorempreskip{0pt}%
\theorempostskip{0pt}%
\newtheorem{proofof}{Proof of\!}%

\newenvironment{proofExt}[1]{{\emph{#1:}}}{\hfill{\hfill\myqedsymbol}}
\newenvironment{proofnoproof}{{}}{{\hfill{\hfill\myqedsymbol}}}

%%%%%%%%%%%%%%%%%%%%%%%%%%%%%%%%%%%%%%%%%%%%%%%%%%%%%%%%%%%%%%%%%%%%
%%%%%%%%%%%%%%%%%%%%%%%%%%%%%%%%%%%%%%%%%%%%%%%%%%%%%%%%%%%%%%%%%%%%
% Macros to embed figures in text

\newlength{\RecordedWidth}
\newcommand{\RecordWidth}[1]{\settowidth{\RecordedWidth}{#1}}%

\newlength{\msize}%

%-------------------------------------------------------------------
% \minipageWFigure{ TEXT }{ Figure }{ Figure caption/label }
%-------------------------------------------------------------------
\newcommand{\minipageWFigure}[3]{%
   \RecordWidth{#2}%
   \setlength{\msize}{\dimexpr(\textwidth-1.1\RecordedWidth)\relax}
   \noindent%
   \SaveIndent%
   \begin{minipage}{\msize}%
       \RestoreIndent%
       #1
   \end{minipage}%
   \null\hfill%
   \begin{minipage}{\RecordedWidth}
       \null\hfill%
       #2%       

       #3%
   \end{minipage}%
}

\newcommand{\EBRY}[2]{#1_{\langle#2\rangle}}%

\newcommand{\NiceFrac}[2]{%
   {\raisebox{0.2em}{\ensuremath{#1}}\!\!\!}%
   {\rotatebox[origin=c]{-20}{\Big/}}% Rotated /
   {\! \raisebox{-0.2em}{\ensuremath{#2}} }}%

% end - Macros to embed figures in text
%%%%%%%%%%%%%%%%%%%%%%%%%%%%%%%%%%%%%%%%%%%%%%%%%%%%%%%%%%%%%%%%%%%% 
%%%%%%%%%%%%%%%%%%%%%%%%%%%%%%%%%%%%%%%%%%%%%%%%%%%%%%%%%%%%%%%%%%%%

%%%%%%%%%%%%%%%%%%%%%%%%%%%%%%%%%%%%%%%%%%%%%%%%%%%%%%%%%%%%%%%%%%
% Color in math
\providecommand{\Mh}[1]{{#1}}%
\providecommand{\IfColorMath}[2]{#2}%

\IfFileExists{.latex_printer_friendly}{\def\GenPrinterVer{1}}{}%

\ifx\GenPrinterVer\undefined
   \IfFileExists{.latex_color}{\def\GenColorMath{1}}{}
\else
   \renewcommand{\IfPrinterVer}[2]{#1}%
\fi

\ifx\GenColorMath\undefined
\else
\renewcommand{\IfColorMath}[2]{#1}%
\renewcommand{\Mh}[1]{{\textcolor{red}{#1}}}%
\fi
% Color in math - end
%%%%%%%%%%%%%%%%%%%%%%%%%%%%%%%%%%%%%%%%%%%%%%%%%%%%%%%%%%%%%%%%%%

\newcommand{\HSetY}[2]{\Mh{H}\pth{#1, #2}}%
\newcommand{\CHSetY}[2]{\Mh{C}\pth{#1, #2}}%
\newcommand{\Holder}{H\"older\xspace}
\newcommand{\OldOm}{{\mathcal{O}}}
\newcommand{\NewOm}{{\mathcal{N}}}

%%%%%%%%%%%%%%%%%%%%%%%%%%%%%%%%%%%%%%%%%%%%%%%%%%%%%%%%%%%%%%%%%%%%%%%
%---------------------------------------------------------------
\newcommand{\IntRange}[1]{\left\llbracket #1 \right\rrbracket}
\newcommand{\IRX}[1]{\left\llbracket #1 \right\rrbracket}%
\newcommand{\IRY}[2]{\left\llbracket #1:#2 \right\rrbracket}

%---------------------------------------------------------------
%%%%%%%%%%%%%%%%%%%%%%%%%%%%%%%%%%%%%%%%%%%%%%%%%%%%%%%%%%%%%%%%%%%%%%%

\newcommand{\obj}{\Mh{f}}%
\newcommand{\Wg}{\Mh{\omega}}%
\newcommand{\TotalW}{\Mh{W}}%

\newcommand{\rankY}[2]{{#1}_{ \permut{#2}}}%%
\newcommand{\areaX}[1]{\Mh{\mathrm{area}}\pth{#1}}%%

\newcommand{\mybinom}[2]%
{\Bigl(\begin{array}{@{}c@{}}#1\\#2\end{array}\Bigr)}

\providecommand{\Set}[2]{\left\{ #1 \;\middle\vert\; #2 \right\}}
\newcommand*{\biga}[1]{\vcenter{\hbox{\scalebox{1.5}{\ensuremath#1}}}}
\newcommand{\CNF}{\TermI{CNF}\xspace}%
\newcommand{\Tree}{\Mh{T}}%

\DeclareSymbolFont{sfletters}{OML}{cmbrm}{m}{it}

\DeclareMathSymbol{\salpha}{\mathord}{sfletters}{"0B}

%%%%%%%%%%%%%%%%%%%%%%%%%%%%%%%%%%%%%%%%%%%%%%%%%%%%%%%%%%%%%%%%%%%%

\providecommand{\ProblemInd}[2]{\textsf{\textcolor{red}{#1}\index{Problem!#2}}}
\newcommand{\NPProblemInd}[4]{
   \begin{center}
      \vspace{0.6cm}
       \fbox{
          \begin{minipage}{0.9\linewidth}
              \vspace{-0.6cm} \hspace{-0.35cm}
              {\large \bfseries \ProblemInd{#1}{#2}}
              \vspace{+0.2cm}%

              \noindent{\textcolor{blue}{\bfseries \textsf{Instance}}:} {#3}
              
              \noindent{\textcolor{blue}{\bfseries \textsf{Question}:}} {#4}
          \end{minipage}
       }
   \end{center}
   \smallskip
   }

\newcommand{\NPProblemOpt}[4]{\NPProblemInd{#1}{#2}{#3}{#4}}

%%%%%%%%%%%%%%%%%%%%%%%%%%%%%%%%%%%%%%%%%%%%%%%%%%%%%%%%%%%%%%%%%%%%

\newcommand{\Hastad}{H{\aa}stad\xspace}

\newcommand{\ANNbr}{\Term{ANN}\index{approximate nearest neighbor}\xspace}
\newcommand{\LSH}{\TermI{L{S}H}\xspace}

\newcommand{\DS}{\EuScript{D}} % DataStructure
\newcommand{\DSC}{\mathsf{D}} % DataStructure

\newcommand{\distH}[2]{\DistChar_{H}\pth{#1, #2}}
\newcommand{\distHY}[2]{\DistChar_{H}\pth{#1, #2}}

\newcommand{\DSAprxNearNbr}{\ensuremath{\DataStructure{{\approx%
         \text{Near{}}}}}%
   \index{near neighbor!data-structure!approximate}%
}
\newcommand{\DataStructure}[1]{\ensuremath{\DS_{#1}}}
\newcommand{\DistNN}[2]{\DistChar\pth{#1,#2}}
\newcommand{\pntc}{\Mh{{p}}}%

\newcommand{\seq}{{\Mh{M}}}%
\newcommand{\cTimes}{\Mh{\beta}}%

\newcommand{\DSTimes}{\Mh{L}}%
\newcommand{\rad}{\Mh{r}}%

\newcommand{\HypercubeX}[1]{\Hypercube^{#1}}
\newcommand{\kFunctions}[2]{\mathrm{combine}\pth{#1,#2}}
\newcommand{\WEqualChar}{\Bumpeq}
\newcommand{\WEqual}[2]{{#1}\WEqualChar{#2}}
\newcommand{\FamilyAmp}{\FamilyC_{\WEqualChar}}
\newcommand{\PSPACE}{\ComplexityClass{PSPACE}\xspace}%

\newcommand{\PS}{\Mh{P}}%

\newcommand{\PSetFar}{\ensuremath{\PntSet^{}_{\geq}\xspace}}%
\newcommand{\DA}{\Mh{D}}%
\newcommand{\uniqX}[1]{\Mh{\mathrm{uniq}}\pth{#1}}%
\newcommand{\diffX}[1]{\mathop{}\!\mathrm{d}#1}

\newcommand{\Frq}{\Mh{F}}%
\newcommand{\frq}{\Mh{f}}%
\newcommand{\STRM}{\Mh{\mathcal{S}}}%
\newcommand{\strm}{\Mh{{s}}}%

\providecommand{\Turan}{Tur\'an\xspace}
\newcommand{\fGrowX}[1]{\Mh{G}^{}_{\Dim}\pth{#1}}
\newcommand{\CompressFunc}{\mathtt{Compress}}
\newcommand{\RSA}{\Term{RSA}\xspace}
\newcommand{\MST}{\Term{MST}\xspace}
\providecommand{\CClassPP}{\CClass{P{}P}\xspace}
\providecommand{\ID}{\mathrm{id}}

\newcommand{\EncodeFunc}{\text{\ttfamily{Enc}}}
\newcommand{\DecodeFunc}{\text{\ttfamily{Dec}}}
\newcommand{\discX}[1]{\mathrm{disc}\pth{#1}}%

\newcommand{\propertyC}{\psi}%
\newcommand{\propertyX}[1]{\propertyC\pth{#1}}%
\newcommand{\numC}{\Gamma}%

\newcommand{\degreeX}[1]{\mathrm{degree}\pth{#1}}

\newcommand{\ind}{m}

\providecommand{\nequiv}{\not\equiv}
\providecommand{\sqcupplus}{\sqcup^+}

\newcommand{\AlgorithmI}[1]{{\protect\color{RedViolet}{#1}}}

\MakePerPage{footnote} % Resets footnote counter on every page

\allowdisplaybreaks % Allows page breaks inside multiline math environments

%-----------------------------------------------------------------------------
% Bibliography & Document Begin
%-----------------------------------------------------------------------------
\IfFileExists{\jobname.bib}{%
   \addbibresource{\jobname.bib}%
}{%
   \addbibresource{shortcuts.bib}%
   \addbibresource{geometry.bib}%
}
\makeindex
\begin{document}
%
\fi
\ChapterNumPage{41}{320}

\Chapter{Double sampling}

\ChapterPrefix{}
\ChapterGen{}
\chaplab{double:sampling}

\remX{Review copy. Mathematical statements and displays retain the
original text; remarks point to verified corrections and shorter proofs
in review.md. Proposed prose additions are blue.}


\QuotePExt{``What does not work when you apply force, would work when
   you apply even more force.''}%
{}%
{Anonymous}

\bigskip%


%%%%%%%%%%%%%%%%%%%%%%%%%%%%%%%%%%%%%%%%%%%%%%%%%%%%%%%%%%%%%%%%%% 
%%%%%%%%%%%%%%%%%%%%%%%%%%%%%%%%%%%%%%%%%%%%%%%%%%%%%%%%%%%%%%%%%% 


% \providecommand{\obj}{f}%
% \providecommand{\Wg}{\omega}%
% \newcommand{\TotalW}{W}%

\section{Double sampling}

Double sampling is the idea that two random independent samples should
look similar, and should \emph{not} be completely different in the way
they intersect a certain set. We use the following sampling model,
which makes the computations somewhat easier.

\begin{defn}
    \deflab{m:sample}%
    %
    Let $\ObjSet = \brc{ \obj_1, \ldots, \obj_n}$ be a set of objects,
    where the $i$\th object has weight $\Wg_i > 0$, for all $i$. Let
    $\TotalW = \sum_i \Wg_i$.  For a target size $\rho$, a
    \emphOnly{$\rho$-sample}\index{$\rho$-sample} is a random sample
    $\Sample \subseteq \ObjSet$, where object $\obj_i$ is picked
    independently with probability
    $\displaystyle \rho \Wg_i / \TotalW$. To simplify the discussion,
    we assume that $\rho \Wg_i / \TotalW < 1$. Handling the more
    general case is easy if somewhat tedious.
\end{defn}

\remX{G3: Specify a positive target size and explicitly require the two
samples below to be independent. See the parameter conventions in review.md.}


\begin{lemma}
    \lemlab{double:sample}%
    %
    Let $\Sample_1$ and $\Sample_2$ be two $\rho$-samples, and
    consider the merged sample $\Sample = \Sample_1 \cup
    \Sample_2$. Let $\SetB \subseteq \ObjSet$ be a set of $m$
    elements. Then, we have that
    \begin{align*}
      \Prob{ \Bigl. \SetB \subseteq \Sample_1 \bigm\vert{ \SetB \subseteq
      \Sample}} &\geq \frac{1}{2^m}
      \quad \text{ and }
      \\
      \Prob{ \Bigl. \SetB \subseteq \Sample_1%
      \text{ and }\;
      \SetB \cap \Sample_2 = \emptyset \bigm\vert{ \SetB \subseteq \Sample}}
      &\leq \frac{1}{2^m}.
    \end{align*}
\end{lemma}
\begin{proof}
    Consider an object $\obj \in \SetB$, and observe that
    \[
        \Prob{\bigl. \obj \in \Sample_1
           \text{ or }  \obj \in \Sample_2
           \mid \obj \in \Sample} = 1.%
    \]
    As such, 
    by symmetry
    \begin{align*}
        % \tau%
        % =%
      \Prob{\bigl. \obj \in \Sample_1 \mid \obj \in \Sample}%
      =%
      \Prob{\bigl. \obj \in \Sample_2 \mid \obj \in \Sample}%
      \geq%
      1/2,
    \end{align*}    
    Now, let $\SetB = \brc{ \obj_1, \ldots, \obj_m}$. Since
    $\Sample_1$ and $\Sample_2$ are independent, and each element is
    being picked independently, we have that
    \begin{align*}
%      \nu%
%      &=%
        \Prob{ \bigl. \SetB \subseteq \Sample_1 \mid{ \SetB \subseteq
        \Sample}}%
        %
      &=%
        \Prob{ \bigl. \obj_1, \ldots, \obj_m  \in \Sample_1 \mid{
        \obj_1, \ldots, \obj_m \in \Sample}}
      \\&%
        =%
        \prod_{i=1}^b \Prob{ \bigl. \obj_i  \in \Sample_1 \mid{
        \obj_1, \ldots, \obj_m \in \Sample}}
      \\&%
          =%
          \prod_{i=1}^b \Prob{ \bigl. \obj_i  \in \Sample_1 \mid{
          \obj_i \in \Sample}}
          \geq%
          \frac{1}{2^m}.
    \end{align*}

    For the second claim, observe that again, by
    symmetry, we have that 
    \begin{equation*}
        \ProbCond{\bigl. \obj \in \Sample_1
           \text{ and }
           \obj \notin \Sample_2
        }{ \obj \in \Sample}%
        =%
%  \end{equation*}
%  \begin{equation*}
      \Prob{\bigl. \obj \notin \Sample_1
      \text{ and } \obj \in \Sample_2
      \mid \obj \in \Sample}%
      \leq%
      1/2,
  \end{equation*}
    as the two events are disjoint. Now, the claim follows by arguing
    as above.
\end{proof}

\subsection{Disagreement between samples on a specific set}

\remX{N2, S1: In the preceding proof, both product limits should be
$m$, not the undefined $b$. The two objectwise conditional probabilities
give a shorter proof, valid even after fixing the entire union.}

We provide \chgY{three}{two} proofs of the following lemma -- the constants are
somewhat different for each version.
\begin{lemma}
    \lemlab{double:sample:2:sets}%
    %
    Let $\Sample_1$ and $\Sample_2$ be two $\rho$-samples from a
    ground set $\ObjSet$, and consider a fixed set $\SetB \subseteq
    \ObjSet$. We have that 
    \begin{equation*}
        \Prob{\Bigl.
           \cardin{\bigl.\cardin{\Sample_1 \cap \SetB} - \cardin{\Sample_2
                 \cap \SetB} }%
           > %
           \eps \rho
        }%
        \leq %
        3 \exp\pth{ -   \eps^2 \rho/2 }.
    \end{equation*}
\end{lemma}

\begin{proof}[Simplest proof.]
    \remX{E1, G1: The stated discrepancy bound is false. A valid
    replacement for $0<\eps\leq1$ is $3\exp(-\eps^2\rho/8)$.
    Both proofs must address the fixed subset, not just the whole
    universe. See the counterexample and complete repair in review.md.}
    By Chernoff's inequality, for $\delta \in (0,1)$, we have
    \begin{equation*}
        \Prob{ \Bigl.\cardin{\bigl.\cardin{\Sample_1} - \rho}
           \geq (\eps/2) \rho } %
        \leq %
        2 \exp\pth{ - (\eps/2)^2 \rho /4 }
        = %
        2 \exp\pth{ - \eps^2 \rho /16 }.
    \end{equation*}
    The same holds for $\Sample_2$, and as such we have
    \begin{align*}
      \Prob{ \Bigl.\cardin{\bigl.\cardin{\Sample_1} - \cardin{\Sample_2}}
      \geq \eps \rho } %
      &\leq %
        \Prob{ \Bigl.\cardin{\bigl.\cardin{\Sample_1} -\rho }
        + \cardin{\bigl.\rho - \cardin{\Sample_2}}
        \geq \eps \rho } %
      \\&
      \leq
      \Prob{ \Bigl.\cardin{\bigl.\cardin{\Sample_1} -\rho }
      \geq (\eps/2)\rho}
      \\&\quad
      +
      \Prob{ \Bigl.
      \cardin{\bigl.\rho - \cardin{\Sample_2}}
      \geq (\eps/2) \rho } %
      \\&
      \leq %
      4 \exp\pth{ - \eps^2 \rho /16 }
    \end{align*}       
\end{proof}

\begin{proof}
    For an object $\obj_i \in \ObjSet$, let $X_i$ be a random 
    variable, where
    \begin{equation*}
        X_i = \begin{cases}
            1 & \obj_i \in \Sample_1 \text{ and } \obj_i \notin
            \Sample_2\\
            -1 & \obj_i \notin \Sample_1 \text{ and } \obj_i \in
            \Sample_2\\
            0 & \text{otherwise.}
        \end{cases}
    \end{equation*}
    We have that
    $p_i = \Prob{X_i = 1} = \Prob{X_i=-1} = (\rho \Wg_i / \TotalW)( 1-
    \rho \Wg_i / \TotalW)$ and $\Ex{X_i} = 0$. Applying the regular
    concentration inequalities in this case is not immediate, since
    there are many $X_i$s that are zero. To overcome this, let $T$ be
    a random variable that is the number of variables in
    $X_1, \ldots, X_n$ that are non-zero. We have that $T$ is a sum of
    $n$ independent $0/1$ random variables, where
    $\Ex{T} = \sum_i 2p_i =2 \rho$. In particular, by
    \lemrefY{Chernoff's inequality}{chernoff:small:delta}, we have
    that
    \begin{equation*}
        q_1 = \Prob{\bigl. T  > (1+\eps) 2 \rho}%
        \leq%
        \exp\pth{ \bigl.- 2\rho\eps^2 /4}
        =%
        \exp\pth{ \bigl.- \rho\eps^2 /2}.
    \end{equation*}
    and assume this happens.  In particular, let $Z_1, \ldots, Z_T$ be
    the non-zero variables in $X_1,\ldots, X_n$, and observe that
    $\Prob{Z_i=1} = \Prob{Z_i=-1} = 1/2$.  Let
    $Y=\sum_i X_i = \sum_i Z_i$. Observe that $\Ex{Y} = 0$, and by
    \thmrefY{Chernoff inequality}{Chernoff:special}, we have that
    \begin{align*}
      q_2%
      &=%
        \Prob{\Bigl.
        \cardin{\bigl.\cardin{\Sample_1 \cap \ObjSet} - \cardin{\Sample_1
        \cap \ObjSet} }%
        > %
        \eps \rho
        }
      \\&
        =%
        \Prob{\Bigl. \cardin{Y - \Ex{Y}} \geq \eps \rho}
      \\&
        \leq%
        \Prob{\Bigl. \cardin{\bigl.\textstyle{\sum}_i Z_i - 0} \geq
        \eps \rho}
      \\%
      &
        \leq%
        2 \exp\pth{ - 2 \frac{ (\eps \rho)^2 }{2T}}
      \\&
        \leq%
        2 \exp\pth{ - 2 \frac{ (\eps \rho)^2 }{2(1+\eps)\rho}} + q_1
      \\&
        =%
        2 \exp\pth{ -  \frac{ \eps^2 \rho }{1+\eps}} + q_1
      \\&
        \leq%
        3 \exp\pth{ -   \eps^2 \rho/2 },
    \end{align*}
    using $T \leq (1+\xi)2\rho$.
\end{proof}



\remX{E2: The second proof has an incorrect mean, missing factors in
the sign bound, and missing conditioning. It also subtracts a sample
from itself and uses an undefined parameter. The fixed set and the
support count both render as $T$. Repair the argument as a unit; see E2.}

\subsection{Exponential decay for a single set}



\begin{lemma}%[Exponential decay for a single set]
    \lemlab{single:decay}%
    %
    Consider a set $\ObjSet$ of $m$ objects, where every object
    $\obj_i \in \ObjSet$ has weight $\Wg_i > 0 $, and
    \begin{math}
        \TotalW= \sum_{i=1}^m \Wg_i.
    \end{math}
    Next, consider a set $\range \subseteq \ObjSet$ such that
    $\WeightX{\range} \geq t \TotalW / \rho$ (such a set is
    \emphOnly{$t$-heavy}\index{$t$-heavy}). Let $\Sample$ be a $\rho$-sample from
    $\ObjSet$. Then, the probability that $\Sample$ misses $\range$ is
    at most $e^{-t}$. Formally, we have
    \begin{math}
        \Prob{ \range \cap \Sample = \emptyset} \leq \exp\pth{-t}.
    \end{math}
\end{lemma}

\begin{proof}
    Let $\range = \brc{\obj_1, \ldots, \obj_k}$.
    \chgY{Clearly, the probability that}{The probability that}
    $\Sample$ \chgY{fails to pick one of these conflicting objects,}{misses the entire range}
    is bounded by
    \begin{math}
      \Prob{ \range \cap \Sample = \emptyset}%
      %&
      =%
      \Prob{ \bigl. \forall i \in \brc{1,\ldots, k} \quad \obj_i
         \notin \Sample_2}%
      =%
      \prod_{i=1}^k \pth{1 - \rho\frac{\Wg_i}{\TotalW}}%
      %
%          \hspace{-2cm}%
      \leq %
      \prod_{i=1}^k \exp \pth{ - \rho \frac{\Wg_i}{\TotalW}} %
        %\\&
        =%
        \exp \pth{ - \frac{\rho}{\TotalW} \sum_i \Wg_i } %
        \leq%
        \exp\pth{ - \frac{\rho}{\TotalW} \cdot t \frac{\TotalW}{\rho}
        } = \exp\pth{-t}. %
      \end{math}
\end{proof}

\subsection{Moments of the sample size}

\remX{N3: In the preceding product calculation, replace the undefined
second sample by the single sample defined in that lemma.}

\newcommand{\cEFunc}[1]{U\pth{#1}}
\newcommand{\cFuncX}[1]{f\pth{#1}}%
\begin{lemma}
    \lemlab{two:samples:gen}%
    %
    Let $\Sample$ be an $m$-sample. \chgY{And let}{Let}
    $\cFuncX{t} \leq \alpha t^\beta$, where $\alpha\geq 1$ and
    $\beta\geq 1$ are constants, such that $m \geq 16\beta$. Then
    \begin{math}
        \cEFunc{m} = \Ex{\bigl.  \cFuncX{\bigl. \cardin{\Sample}} } \leq 2
        \alpha (2m)^\beta.
    \end{math}
\end{lemma}
\begin{proof}
    The proof follows from Chernoff's inequality and some tedious but
    straightforward calculations. The reader is as such encouraged to
    skip reading it.

    Let $X = \cardin{\Sample}$. This is a sum of $0/1$ random variables with
    expectation $m$. As such, we have
    \begin{equation*}
        \nu%
        =%
        \Ex{\bigl. \cFuncX{\bigl.\cardin{\Sample}} }%
        \leq %
        \sum_{i=0}^\infty \Prob{X = i} f(i)
        \leq%
        \alpha 
        \sum_{i=0}^\infty \Prob{X = i} 
        i^\beta.        
    \end{equation*}
    \remX{E3, S2: The following binning chain is false. Split directly
    at $2m$ to repair it. Alternatively, a short exponential-moment
    argument yields $\alpha(2m)^\beta$ under weaker
    hypotheses; see review.md.}
    Considering the last sum, we have
    \begin{align*}
        \sum_{i=0}^\infty \Prob{\big. X = i} i^\beta
        &\leq %
        \sum_{j=0}^\infty \Prob{\big. X \geq  j  m } \pth{(j+1) m}^\beta%
        \\&
        \leq %
        (2m)^\beta
        + m^\beta \sum_{j=2}^\infty \Prob{\big. X \geq  j  m }
        \pth{j+1 }^\beta.%
    \end{align*}
    By Chernoff's inequality (see \thmref{Chernoff:simplified}), we
    can bound the last summation as
    \begin{align*}
      \tau%
      &= %
        \sum_{j=2}^5 \Prob{\big. X \geq  j  m }
        \pth{j+1 }^\beta
        + 
        \sum_{j=6}^\infty \Prob{\big. X \geq  j  m }
        \pth{j+1 }^\beta
      \\&%
      \leq      
      \sum_{j=2}^5 \exp \pth{ - \frac{m (j-1)^2}{4} }
      \pth{j+1 }^\beta
      + 
      \sum_{j=6}^\infty 2^{-jm}
      \pth{j+1 }^\beta %
      \\&%
      \leq%
      \exp\pth{ - \frac{m}{4} } 3^\beta
      + \exp\pth{ - m} 4^\beta
      + \exp\pth{ - 2m} 5^\beta
      + \exp\pth{ - 4m} 6^\beta
      \\&\quad
      + \sum_{j=6}^\infty 2^{-jm}
      \pth{j+1 }^\beta %
      < 1,
    \end{align*}
    since $m \geq 16 \beta$. We conclude that
    $\nu \leq \alpha (2 m)^\beta + \alpha m^\beta \tau \leq 2
    \alpha(2m)^\beta$.
\end{proof}

\begin{remark}
    The constant $16$ in the above lemma is somewhat strange. A better
    constant can be derived by breaking the range of sizes into
    smaller intervals and using the right Chernoff inequality. Since
    this is somewhat tangential to the point of this write-up, we
    leave it as is (i.e., this constant is not critical to our
    discussion).
\end{remark}



\subsection{Growth function}

\remX{E4: Define the extremal function using VC dimension at most
$\Dim$. The last inequality below is false already for $\Dim=1$:
the sum is $n+1$. Delete that last inequality.}

The \emphOnly{growth function}\index{growth function} $\Growth{\Dim}{n}$ is the maximum number
of ranges in a range space with \VC dimension $\Dim$, and with $n$
elements. By Sauer's lemma, it is known that 
\begin{equation}
    \Growth{\Dim}{n} = \sum_{i=0}^{\Dim} \binom{n}{i}
    \leq \sum_{i=0}^{\Dim} \frac{n^i}{i!}
    \leq n^\Dim,
    % 
    \eqlab{g:d:n:estimate:2}
\end{equation}

\chgY{The following is well known (the estimates are somewhat
tedious to prove)}{The following estimates appear in Har-Peled}~\cite{h-gaa-11}:
\remX{S3: A generating-function argument proves a stronger upper
bound without the factor two; see review.md.}
\begin{lemma}[\cite{h-gaa-11}]
    \lemlab{g:estimate}%
    %
    For $n \geq 2\Dim$ and $\Dim \geq 1$, we have
    \begin{math}
        \displaystyle \pth{\frac{n }{\Dim}}^\Dim%
        \leq%
        \Growth{\Dim}{n}%
        \leq%
        2 \pth{\frac{n e}{\Dim}}^\Dim,
    \end{math}
    where $\Growth{\Dim}{n} = \sum_{i=0}^{\Dim} \binom{n}{i}$.
\end{lemma}

\newcommand{\Net}{{\Sample}}%
\newcommand{\NetA}{{T}}%

\remX{G2, S3: The next proof applies estimates outside their stated
hypotheses and fails on small realized sample sizes. A factorial-moment
argument proves the stronger bound $(2em/\Dim)^\Dim$ for
$\Dim\geq1$ and $m\geq\Dim$, bypassing the preceding moment lemma.
The duplicated sample also needs a distinct symbol.}

\begin{lemma}
    \lemlab{two:samples}%
    %
    Let $\Sample$ and $\SampleA$ be two independent $m$-samples from
    $\FGroundSet$. Assume that $m \geq \Dim$. Then
    \begin{math}
        \Ex{\bigl. \Growth{\Dim}{\bigl.\cardin{\Sample} +
              \cardin{\SampleA}} }%
        \leq%
        \fGrowX{2m},
    \end{math}
    where 
    \begin{math}
        \fGrowX{2m} = 4 \, \pth{{4 e m}/{\Dim}}^{\Dim}.
  \end{math}
\end{lemma}
\begin{proof}
    We set
    \begin{math}
        \alpha = 2 \pth{\frac{ e}{\Dim}}^\Dim,
    \end{math}
    $\beta = \Dim$, and $f(n) = \alpha n^\beta$. Duplicate every
    element in $\FGroundSet$, and let $\FGroundSet'$ be the resulting
    set. Clearly, the size of a $2m$-sample $\Sample$ from
    $\FGroundSet'$ is the same as $\cardin{\Sample} + \cardin{T}$.  By
    \lemref{g:estimate}, we have
    \begin{math}
      \Ex{\Bigl. \Growth{\Dim}{\bigl.\cardin{\Sample}} }%
      \leq %
      \Ex{\Bigl. f \pth{\bigl.\cardin{\Sample}} }%
      \leq %
      2 \alpha (4m)^\beta%
      \leq%
      4 \pth{\frac{ 4em}{\Dim}}^\Dim
    \end{math}
    The last inequality follows from \lemref{two:samples:gen}.
\end{proof}



%%%%%%%%%%%%%%%%%%%%%%%%%%%%%%%%%%%%%%%%%%%%%%%%%%%%%%%%%%%%%%%%%%%%%%%%%
%%%%%%%%%%%%%%%%%%%%%%%%%%%%%%%%%%%%%%%%%%%%%%%%%%%%%%%%%%%%%%%%%%%%%%%%%
%%%%%%%%%%%%%%%%%%%%%%%%%%%%%%%%%%%%%%%%%%%%%%%%%%%%%%%%%%%%%%%%%%%%%%%%%

\section{Proof of the $\eps$-net theorem}
\seclab{full:proof}

\chgY{Here we are working in the unweighted settings (i.e., the weight of a
single element is one)}{Here every element has weight one}.
\chgY{}{The theorem is due to Haussler and Welzl}~\cite{hw-ensrq-87}.
\remX{G3: Require $0<\BadProb<1$. Handle dimension zero separately
before dividing by $\Dim$. If the requested size reaches the ground-set
size, take the full ground set.}

\begin{theorem}[$\eps$-net theorem, \cite{hw-ensrq-87}]
    \index{random sample}%
    \index{net!theorem@$\varepsilon$-net theorem|textbf}%
    \index{theorem!net@$\varepsilon$-net|textbf}%%
    \thmlab{epsilon:net:2}%
    % 
    Let $\Dim$ be the \VC dimension of a range space
    $(\GroundSet,\RangeSet)$, let $\FGroundSet$ be a finite subset of $\GroundSet$, and
    suppose that $0 < \eps \leq 1$ and $\BadProb < 1$. Let $N$ be an
    $m$-sample from $\FGroundSet$ (see \defref{m:sample}), where
    \begin{equation}
        m \geq \max\pth{\;
           \frac{8}{\eps} \lg \frac{4}{\BadProb} \,,\,
           \frac{16\Dim}{\eps}\lg \frac{16}{\eps} \;}.
        \eqlab{e:net:2}
    \end{equation}
    Then $N$ is an $\eps$-net for $\FGroundSet$ with probability at
    least $1-\BadProb$.
\end{theorem}


\subsection{The proof}

\subsubsection{Reduction to double sampling}

Let $n = \cardin{\FGroundSet}$.  Let $N$ be the $m$-sample from
$\FGroundSet$.  Let $\Event_1$ be the \chgY{probability}{event} that $N$ fails to be
an $\eps$-net.  Namely,
\begin{equation*}
    \Event_1%
    =%
    \Set{ \exists{\range \in \RangeSet} }{\bigl.\ts
       \cardin{\range \cap \FGroundSet} \geq \eps n \text{ ~ and ~
       } \range \cap N = \emptyset}.    
\end{equation*}
(Namely, there exists a ``heavy'' range $\range$ that does not contain
any point of $N$.)  To complete the proof, we must show that
$\Prob{\Event_1}\leq \BadProb$. Let $T$ be another $m$-sample
generated in a similar fashion to $N$. Let $\Event_2$ be the event
that $N$ fails but $T$ ``works''. Formally
\begin{equation*}
    \Event_2%
    =%
    \Set{\bigl. \exists{\range \in \RangeSet}}%
    {\bigl.\cardin{\range \cap \FGroundSet} \geq \eps n, \, \range \cap
       N = \emptyset, \, \text{ and } \cardin{\range \cap T} \geq
       \frac{\eps m}{2} }.
\end{equation*}
Intuitively, since
$\Ex{\bigl. \cardin{\range \cap \FGroundSet} } \geq \eps m$, we have
that for the range $\range$ that $N$ fails for, it follows with
``good'' probability that $\cardin{\range \cap T} \geq \eps m /
2$. Namely, $\Event_1$ and $\Event_2$ have more or less the same
probability.

\remX{E5: In the preceding intuition, the random intersection is with
$T$, not the ground set. In the following proof, condition on $N$,
choose its heavy witness, and draw $T$ independently. Correct the
$n/m$ threshold mismatch, strict-event comparison, and exponential
equality; see E5 for the complete argument.}

\begin{claim}
    \clmlab{V:C:e1}%
    %
    $\Prob{\Event_2} \leq \Prob{\Event_1} \leq 2\Prob{\Event_2}$.
\end{claim}

\begin{proof}
    Clearly, $\Event_2 \subseteq \Event_1$, and thus $\Prob{\Event_2}
    \leq \Prob{\Event_1}$. As for the other part, note that by the
    definition of conditional probability, we have
    \[
        \Prob{ \Event_2 \mid{ \Event_1}} = \Prob{\Event_2 \cap
           \Event_1} / \Prob{\Event_1} = \Prob{\Event_2} /
        \Prob{\Event_1}.
    \]
    It is thus enough to show that $\Prob{\Event_2 \mid{\Event_1}}$
    $\geq 1/2$.
    
    Assume that $\Event_1$ occurs. There is $\range \in \RangeSet$,
    such that $\cardin{\range \cap \FGroundSet} > \eps n$ and
    $\range \cap N = \emptyset$.  The required probability is at least
    the probability that for this specific $\range$, we have
    $X = \cardin{ \range \cap T} \geq \tfrac{\eps n}{2}$. The variable
    $X$ is a sum of $t = \cardin{\range \cap \FGroundSet} \geq \eps n$
    random independent 0/1 variables, each one has probability $m/n$
    to be one. Setting $\mu = \Ex{X} = t m/n \geq \eps m$ and
    $\xi = 1/2$, we have by Chernoff's inequality that
    \begin{align*}
      \Prob{ \cardin{\range \cap T} \leq \eps m /2} %
      \leq%
      \Prob{ X  < (1-\xi) \mu }%
      \\&
      \leq%
      \exp \pth{ - \mu \xi^2 / 2 }%
      \\&
      =%
      \exp \pth{ - \eps m / 8 }%
      < %
      1/2,
    \end{align*}
    if $\eps m \geq 8$.  Thus, for $\range \in \Event_1$, we have
    \begin{math}
        \Prob{\Event_2} / \Prob{\Event_1} \geq \Prob{\cardin{ \range
              \cap T} \geq \tfrac{\eps m}{2}} = 1 -\Prob{
           \cardin{\range \cap T} < \frac{\eps m}{2} } \geq
        \frac{1}{2}.
    \end{math}
%    \aftermathA
\end{proof}


\clmref{V:C:e1} implies that to bound the probability of $\Event_1$,
it is enough to bound the probability of $\Event_2$.  Let
\begin{equation*}
    \Event_2'%
    =%
    \Set{ \exists{\range \in \RangeSet} }{\bigl. \range \cap
       N
       = \emptyset \text{ and }\cardin{\range \cap T} \geq
       \frac{\eps m}{2} }.
\end{equation*}
Clearly, $\Event_2 \subseteq \Event_2'$. Thus, bounding the
probability of $\Event_2'$ is enough to prove
\thmref{epsilon:net:2}. Note, however, that a shocking thing happened!
We no longer have $\FGroundSet$ participating in our event. Namely, we
turned bounding an event that depends on a global quantity (i.e., the
ground set $\FGroundSet$) into bounding a quantity that depends only
on a local quantity/experiment (involving only $N$ and $T$).  This is
the crucial idea in this proof.

\subsubsection{Using double sampling to finish the proof}


\remX{N1, G4: The macro used for the range trace renders as $T$, the
same symbol as the second sample; use $B$ for the trace. Condition on
the entire fixed union and justify objectwise independence. There are
at most the stated number of traces, and a small trace only rules out
that particular witness, not the whole event.}

\begin{claim}
    \clmlab{e2}%
    %
    \begin{math}
        \Prob{\bigl. \Event_2} \leq \Prob{ \bigl. \Event_2'}%
        \leq%
        2^{-\eps m /2} \fGrowX{2m}.
    \end{math}
\end{claim}
\begin{proof}
    We fix the content of $\Sample = N \cup T$. The range space
    $(\Sample, \VCProj{\RangeSet}{\Sample})$ has
    $\Growth{\Dim}{ \cardin{\Sample}}$ ranges. Fix a range $\range$ in
    this range space. Let $\SetB = \range \cap \Sample$.  If
    $b = \cardin{ \SetB } < \eps m /2$ then the $\Event_2'$ can not
    happened. Otherwise, the probability that $\range$ is a bad range
    is
    \begin{math}
        \Prob{ \Bigl. \SetB \subseteq T \text{ and }\; \SetB \cap N =
           \emptyset \bigm\vert{ \SetB \subseteq \Sample}} \leq
        \frac{1}{2^b},
    \end{math}
    by \lemref{double:sample}. In particular, by the union bound over
    all ranges, we have
    \begin{math}
        \Prob{\Event_2' \mid \Sample} %\Growth{\Dim}{ \cardin{\Sample}}
        \leq 2^{-\eps m /2}  \Growth{\Dim}{ \cardin{\Sample}}.
    \end{math}
    As such, we have
    \begin{align*}
      \Prob{\Event_2'}%
      =%
      \sum_{\Sample}       
      \Prob{\bigl. \Event_2' \mid \Sample} \Prob{\Sample}
      \\&
      \leq%
      \sum_{\Sample}
      2^{-\eps m /2}  \Growth{\Dim}{\bigl. \cardin{\Sample}}
      \Prob{\Sample}
      \\&
      \leq%
      2^{-\eps m /2}
      \Ex{\Bigl.
      \Growth{\Dim}{\bigl. \cardin{\Sample}} }
      \\&
      \leq%
      2^{-\eps m /2}
      \fGrowX{2m}.
    \end{align*}
    by \lemref{two:samples}.
\end{proof}

% \newpage

\begin{proofof} {\textsc{\thmref{epsilon:net:2}.}}
    \remX{E6: The prefactor after substitution is eight, not sixteen.
    Sixteen can be used as a sufficient upper bound. The first bracket
    needs $\lg(4e)$, not $\lg(4e/\Dim)$, and the comparison with
    $-8\Dim$ has its direction reversed. Correcting these steps
    preserves the theorem's sample-size constants; see review.md.}
    By \clmref{V:C:e1} and \clmref{e2}, we have that
    \begin{math}
        \Prob{\Event_1} \leq%
        2 \cdot 2^{-\eps m /2} \fGrowX{2m}.
    \end{math}
    It thus remains to verify that if $m$ satisfies \Eqref{e:net:2},
    then the above is smaller than $\BadProb$.
    Which is equivalent to
    \begin{align*}
      2 \cdot 2^{-\eps m/2} \fGrowX{2m} \leq \BadProb
      &  \iff
      %
        16 \cdot 2^{-\eps m/2} \pth{\frac{4 e m}{\Dim}}^{\!\Dim}
        \leq \BadProb
      \\&
      \iff
      %
        -4 + \frac{\eps m }{ 2} - \Dim \lg \pth{\frac{4 e m}{\Dim}}
        \geq \lg \frac{1}{\BadProb}%
      \\%
      &%
        \iff
        \pth{ \frac{ \eps m}{8} -4 - \Dim \lg \frac{4 e}{\Dim}}
        + \pth{ \frac{\eps m}{8} - \lg \frac{1}{\BadProb}}
        \\&\qquad
        + \pth{\frac{\eps m }{ 4} - \Dim \lg \pth{\frac{ m}{\Dim}}}
        \geq 0%
    \end{align*}
    We remind the reader that the value of $m$ we pick is such that 
    \begin{math}
        m \geq \max\pth{ \frac{8}{\eps} \lg \frac{4}{\BadProb} \,,\,
           \frac{16\Dim}{\eps}\lg \frac{16}{\eps} }.
    \end{math}
    In particular, $m \geq 64 \Dim / \eps$
    and
    \begin{math}
        -4 - \Dim \lg \pth{\frac{4 e }{\Dim}} \geq -4 - 4 \Dim \leq -8
        \Dim \geq -\eps m /8.
    \end{math} 
    Similarly, by the choice of $m$, we have
    \begin{math}
        \eps m /8 \geq \lg \frac{1}{\BadProb}. 
    \end{math}
    As such, we need to show that
    \begin{math}
        \frac{\eps m }{ 4} \geq \Dim \lg \pth{\frac{ m}{\Dim}} \iff%
        m \geq \frac{4\Dim}{\eps} \lg {\frac{ m}{\Dim}},
    \end{math}
    and one can verify using some easy but tedious calculations that
    this holds if $ m \geq \frac{16\Dim}{\eps}\lg \frac{16}{\eps} $.
%    
    % This completes the proof of the theorem.
\end{proofof}







\section{Bibliographical notes}

\remX{W3: The following subsection is the requested addition on
connections beyond range spaces. References are supplied in the local
review bibliography.}

\subsection{Double sampling beyond range spaces}

\chgY{}{Double sampling uses symmetrization: it exploits exchangeability
between independent copies. Vapnik and
Chervonenkis (1971)}~\cite{vc-ucrfe-71}
\chgY{}{used this method to prove uniform convergence of empirical
frequencies. Haussler and Welzl (1987)}~\cite{hw-ensrq-87}
\chgY{}{used these ideas for epsilon-nets and geometric range searching.
The second sample, often called a ghost sample, is needed only for
analysis. Conditioning on the combined sample reduces the collection
of events to finitely many traces, making a union bound effective.}

\chgY{}{In independent holdout validation, a predictor chosen from
training data is evaluated on fresh data from the same distribution.
Conditional on the training data
and the training procedure's randomness, the predictor is fixed.
For bounded loss, independent identically distributed holdout
examples therefore give a concentration bound on its prediction
error. Here the second sample is observed. Adapting the predictor to
holdout results invalidates this argument.
Dwork \etal (2015)}~\cite{dfhprr-rh-15}
\chgY{}{developed mechanisms that control the effects of adaptive
holdout reuse.}

\chgY{}{The variance method of Efron and Stein
(1981)}~\cite{es-jev-81}
\chgY{}{compares a function of independent inputs before and after
replacing one coordinate by an independent copy. This is a related
resampling argument. For a sum of independent random
matrices, subtracting independent copies of the summands produces
independent symmetric differences. Multiplying these differences by
independent uniform random signs preserves their joint distribution.
Such symmetrization is used in sparse signal reconstruction from
Fourier measurements by Rudelson and Vershynin
(2008)}~\cite{rv-srfgm-08}
\chgY{}{and in low-rank matrix completion by Cand\`es and Recht
(2009)}~\cite{cr-emcco-09}.
\chgY{}{These applications recover signals or structured matrices
from incomplete observations.}

\chgY{}{A compact algorithmic example is the single-sample
prophet inequality of Rubinstein \etal
(2020)}~\cite{rww-oscp-20}.
\chgY{}{Independent nonnegative values arrive sequentially, and one
value may be accepted irrevocably. Beforehand, draw one independent
sample from each value's distribution. Assign independent uniform
labels from the unit interval to all samples and arriving values,
and compare value--label pairs lexicographically. Set the threshold
to the largest sample pair and accept the first arriving pair above
it, receiving zero if none qualifies. The expected reward is at least one-half
of the expected maximum of the arriving values, assuming that maximum
has finite expectation. The proof draws two values per distribution
and independently flips fair coins to assign their sample and arrival
roles.}

\StandAloneMode{
   \section{From previous lectures}

   \remX{N5: In the included simplified Chernoff theorem, the failure
   probability should be $1-p_i$, not $1-p$. The shared original
   fragment is left unchanged.}
   
%   \IncFragment{convex:hull}
%   \IncFragment{Chebychev}
   \IncFragment{chernoff_small_delta}
   \IncFragment{Chernoff_simplified}
   % \IncFragment{Chernoff_below}
%   \IncFragment{t:Chernoff:special}
   \IncFragment{c_special}%
 %  \IncFragment{epsilon:sample:stronger
%
}



\ChapterEnd{}

% --------------------------------------------------------
% 
% 12_VC dim - end of file
% --------------------------------------------------------
  
